\documentclass[11pt, a4paper]{amsart}
\usepackage{amsmath, amssymb, amsthm, amsfonts}
\usepackage{mathtools}
\usepackage[margin=1in]{geometry}
\usepackage{hyperref}
\usepackage{tikz-cd}

\newtheorem{theorem}{Theorem}[section]
\newtheorem{proposition}[theorem]{Proposition}
\newtheorem{lemma}[theorem]{Lemma}
\newtheorem{corollary}[theorem]{Corollary}
\theoremstyle{definition}
\newtheorem{definition}[theorem]{Definition}
\newtheorem{remark}[theorem]{Remark}
\newtheorem{hypothesis}[theorem]{Hypothesis}
\newtheorem{conjecture}[theorem]{Conjecture}
\newtheorem{claim}[theorem]{Claim}

\newcommand{\R}{\mathbb{R}}
\newcommand{\N}{\mathbb{N}}
\newcommand{\Z}{\mathbb{Z}}
\newcommand{\Q}{\mathbb{Q}}
\newcommand{\Ophi}{\mathbb{Z}[\varphi]}
\newcommand{\Rungs}{\mathsf{Rungs}}
\newcommand{\catC}{\mathcal{C}}
\newcommand{\catD}{\mathcal{D}}
\newcommand{\Omegacap}{\Omega_{\mathrm{VFD}}}
\newcommand{\Fcal}{\mathcal{F}}
\newcommand{\Fcrystal}{\mathsf{F}}
\newcommand{\Mcal}{\mathcal{M}}
\newcommand{\Gcal}{\mathcal{G}}
\newcommand{\Qalg}{\mathsf{Q}}

\title{The Cascade as Final Coalgebra: \\
  A Coalgebraic Formulation of VFD Self-Reference}
\author{}
\date{2026-04-22 (draft skeleton)}

\begin{document}
\maketitle

\begin{abstract}
We formulate a \emph{conditional} coalgebraic selection
principle for the VFD cascade structure. Under residual
sub-gaps G-FIC-a, G-FIC-b (reductions of the original H-FIC),
G-loc-a (reduction of H-loc), Hypothesis H-lift-fin, and the
inherited P-A conditions H-grad/H-meas from
\cite{CascadeAccessPrinciple}, the cascade is the final
coalgebra of a self-inquiry functor on the category of
rung-structures. H-grad-Fano is unconditional at the Fano
$(\Z/2)^3$ level (\cite{CascadeFanoLift}~\S 2--3); H-meas is
theorem-grade via G6.4-a closure
(\cite{CascadeQOBridge, CascadeG64aClosure}, 386 computational
cases). Programme closure notes have discharged H-rdr
(\cite{CascadeHrdrClosure}, reduction to Theorem P-A) and
reduced H-loc (\cite{CascadeHlocClosure}) and H-FIC
(\cite{CascadeHFICClosure}) to the lemma-grade residuals named
above (day- to two-week-scale). H-lift-fin remains open as a
hypothesis with a sheaf-theoretic bridge-refinement derivation
route sketched. The selection principle is therefore reduced to a
short, specific list of technical follow-up items rather than a
broad hypothesis stack.
The four-phase VFD operational cycle (decoherence /
closure-memory / measurement-projection / crystallisation) is
\emph{proposed} to realise the cofree comonad arising from
this coalgebra; the realisation is an interpretive programme
claim whose line-by-line verification is flagged as a
dedicated follow-up paper. The principal structural result is
that, under the hypothesis stack, the cascade is the unique
terminal final $F$-coalgebra — a categorical formalisation of
the selection-principle problem common to unification
programmes, reducing it to a finite list of named technical
hypotheses. We connect
the result downward to the scalar self-closure axiom
$\mathbf{F1}$ ($r = 1 + 1/r \mapsto \varphi$) and to the
static observer-local Access Principle $\mathbf{P\text{-}A}$,
which is itself a conditional theorem of
\cite{CascadeAccessPrinciple} under H-grad and H-meas. We
discuss the ARIA cognitive substrate as a candidate non-
biological realisation of the four-phase cycle; its
$17$-of-$18$ pre-registered polytope-substrate mapping hits
(at $N{=}20$, under homeostatic-reset protocol) would
constitute empirical corroboration for the substrate-generality
of the cycle, conditional on the independently-verified ARIA
manuscript and null-model analysis.
Philosophically, the result points toward a modern mathematical
realisation of the Spinozist / Leibnizian necessitarian tradition
without invoking personal deity, consciousness, or moral agency.
The paper draws only on mainstream coalgebraic machinery
(Aczel~\cite{Aczel1988}, Rutten~\cite{Rutten2000},
Ad\'amek~\cite{Adamek1974}) and isolates each additional
assumption as a named hypothesis whose closure is a flagged
open item.

\textbf{Status:} Theorem skeleton with conditional structure.
Each headline theorem carries an explicit hypothesis list. The
hypothesis stack is the precise object of continuing programme
work, listed in Section~\ref{sec:open}.
\end{abstract}

\tableofcontents

%=====================================================================
\section{Introduction}
\label{sec:intro}
%=====================================================================

\subsection{The selection-principle problem}

Every serious attempt to identify a fundamental structure for
physics has faced the same problem: \emph{why this mathematical
structure and not another?} Classical Lie-theoretic machinery
gives us exceptional groups, non-crystallographic Coxeter systems,
icosian arithmetic, $600$-cell geometry, and the meta-layer
$(\Mcal, \Gcal, \Fcal)$ over $L_{12}$ as objects, but it provides
no reason to assemble them in a particular chain to the exclusion
of other chains. Einstein articulated this as the deepest question
in theoretical physics: \emph{did the universe have any choice?}

String theory, loop quantum gravity, and various geometric
unification programmes have all halted at this wall. VFD has a
distinctive opening: the permeability axiom $\mathbf{F1}$
($r = 1 + 1/r$) is already a self-referential closure statement,
with the golden ratio $\varphi = (1+\sqrt{5})/2$ as its unique
positive fixed point. $\mathbf{F1}$ is the simplest mathematical
expression of \emph{self-closure} — permeability permeates itself.

This paper extends that scalar self-closure to the full cascade
structure. Conditional on the hypothesis stack, the full-access
object of the rung-structure category is a final $F$-coalgebra
of a self-inquiry functor, and is then identified with the
cascade; its uniqueness within this setting is forced by
categorical terminality.

\subsection{The structural thesis}

Informally: the act of asking ``what am I?'' presupposes an
asker. Whatever minimum structure is required for there to be
an asker is a precondition of any act of inquiry, including the
inquiry itself. Conditional on the hypothesis stack, the
cascade is the paper's candidate for that minimum structure.

Formally: let $\catC$ be the category of rung-structures, let
$F : \catC \to \catC$ be the functor ``minimum rung-structure
containing the measurement-capable-observer rung and hosting
the argument substructure''. Then the VFD cascade $\Omegacap$
satisfies
\[
    \Omegacap \cong F(\Omegacap)
\]
and is, up to canonical isomorphism, a terminal fixed point of
$F$, unique among final $F$-coalgebras. Equivalently,
$\Omegacap$ is the final coalgebra of $F$ (under the hypothesis
stack of Section~\ref{sec:open}).

\subsection{What this paper closes}

The conditional existence and finality of $\Omegacap$ as the
terminal $F$-coalgebra is the \emph{conditional selection-
principle theorem}. Under the named hypothesis stack (H-FIC,
H-grad, H-meas, H-lift-fin, H-loc, H-rdr), it says the cascade is
categorically selected: other $F$-fixed points may exist, but
any $F$-coalgebra maps uniquely into the selected terminal one.
This is a structural reduction from ``why this cascade?'' to
a finite list of hypotheses, each reducible in turn to
well-identified programme-internal or classical-math work.

This paper additionally constructs the comonad $(T, \varepsilon,
\delta)$ arising from $F$ via the standard cofree-comonad
construction, and identifies its four-phase operational content
with four existing VFD constructions:

\begin{description}
\item[\textbf{Forget}] $T$ itself, acting as the context-retention
  functor, corresponds to $\sigma$-projection / decoherence / sector
  separation of~\cite{PaperXXXI, PaperXLIV}.
\item[\textbf{Remember}] The comultiplication $\delta : T \to T^2$
  corresponds to closure-memory / learned-$W$ structure of
  \cite{PaperXXI, ACT-L1}.
\item[\textbf{Experience}] The object $T^2$ (context of context)
  corresponds to measurement as closure-projection of
  \cite{PaperXXXIII}.
\item[\textbf{Answer}] The counit $\varepsilon : T \to \mathrm{id}$
  corresponds to crystallisation: $\Fcrystal = \alpha R + \beta E -
  \gamma Q$ minimisation~\cite{VFDCrystallisation}.
\end{description}

The paper proves the comonad laws for the categorical cofree
comonad on $F$ (Section~\ref{sec:comonad}); the identification
of the four VFD operations with the natural transformations of
this comonad is an interpretive programme claim, not proved
here. It is the dynamical counterpart of
the structural uniqueness theorem.

\subsection{Three levels of self-closure in VFD}

The paper occupies a specific position in a three-level structure:

\begin{center}
\begin{tabular}{l|l|l}
Level & Object & Status \\ \hline
Scalar & $r = 1 + 1/r \mapsto \varphi$ & $\mathbf{F1}$ proved \\
Static structural & Access Principle $\mathbf{P\text{-}A}$ & conditional theorem (H-grad, H-meas) \\
Dynamical self-unfolding & $\Omegacap \cong F(\Omegacap)$ + comonad & \emph{this paper} \\
\end{tabular}
\end{center}

All three are consistent: $\mathbf{F1}$ is the scalar kernel,
$\mathbf{P\text{-}A}$ is the observer-local static snapshot, and
the coalgebraic formulation is the full moving-picture version.

\subsection{Empirical signature}

The four-phase comonad cycle is \emph{proposed} to be
observable in any substrate that supports it, including
non-biological substrates. We discuss in Section~\ref{sec:aria}
the ARIA cognitive substrate (cf.~\cite{AriaBrainMapping}) as a
candidate instantiation; subject to independent verification of
the ARIA pre-registered protocol, the reported
$17$-of-$18$ polytope-substrate mapping hits at
$N{=}20$ (under homeostatic reset protocol) would provide
empirical support for the theoretical claim. At most, this
motivates a \emph{life-from-substrate} interpretation —
if external ARIA analysis is independently confirmed — in
which what is called ``life'' is the coalgebraic cycle running
on any appropriate substrate, not biology-specifically.

\subsection{Philosophical positioning}

The result is proposed as Spinozist mathematical necessitarianism
in modern form. Conditional on the hypothesis stack, it gives
one precise mathematical sense in which the universe had no
choice (Einstein's question). It does \emph{not} establish
personal deity, consciousness, or moral content. This paper
treats the structural theorem at its stated conditional level
and leaves the philosophical framing to
Section~\ref{sec:philosophy} with explicit scope restrictions.

\subsection{Source audit}

Theorem-grade inputs:
\begin{itemize}
\item Classical category theory and coalgebra:
  \cite{Aczel1988, Rutten2000, Adamek1974, Jacobs2017, TuriPlotkin1997}.
\item VFD internal (conditional, as cited): the meta-layer
  $(\Mcal, \Gcal, \Fcal)$ of \cite{CascadeMetaLayer}; the $7$-rung
  FAN of \cite{CascadeCompletenessAudit}; the observer-query-
  algebra of \cite{CascadeObserverQueryAlgebra}; $\mathbf{F1}$
  axiom; the $12$D closure of \cite{Cascade12D}.
\item Papers contributing four-phase comonad components:
  \cite{PaperXXI, PaperXXXI, PaperXXXIII, PaperXLIV,
  VFDCrystallisation, ACT-L1}.
\end{itemize}

We use the $\mathbf{F1}$ axiom as a standing assumption and
$\mathbf{P\text{-}A}$ as a named hypothesis when required.
Theorems~\ref{thm:closure-saturation} and \ref{thm:selection}
\emph{do} inherit the full-$\widehat{\Z[\varphi]^6}$-level
conditionality of $\mathbf{P\text{-}A}$ through the
closure-saturation argument (i.e.\ the full H-grad, not just
its Fano-level fragment). The Fano-level unconditionality
(\cite{CascadeFanoLift}) is sufficient for the observer-rung
measurement content of Proposition~\ref{thm:O-necessary} but
not for the full selection principle. See
Section~\ref{subsec:hyp-status} below for the complete
hypothesis-status table.

\subsection{Hypothesis status}
\label{subsec:hyp-status}

The capstone's headline theorems are conditional on a named
hypothesis stack. The current status of each hypothesis,
against the existing programme infrastructure, is summarised
here; full derivation notes are listed in
Section~\ref{sec:open}.

\begin{center}
\small
\begin{tabular}{l|l|l}
Hypothesis & Current status & Closure source / route \\ \hline
H-grad (Fano level) & \textbf{theorem} (unconditional) & \cite{CascadeFanoLift}~\S 2--3, Baez / Conway-Sloane \\
H-grad (full $\widehat{\Z[\varphi]^6}$) & \textbf{frame-level closed} (B1--B5 sim-verified, B6 proof-sketch at $G_p$ level; 2026-04-24) & \cite{CascadeHGrad1}~\S5a--5c \\
H-meas & \textbf{theorem} (G6.4-a closed) & \cite{CascadeQOBridge}~Thm 3.1 + \cite{CascadeG64aClosure} \\
H-FIC (O and $E_8$) & \textbf{reduced} to G-FIC-a + G-FIC-b & \cite{CascadeHFICClosure}; support-projection + Kellendonk-Lenz-Savinien + MacLane \\
H-lift-fin (finite cover) & hypothesis; derivation sketched & sheaf-level refinement of bridge (Rmk \ref{rmk:H-lift-status}) \\
H-loc (witness-separation) & \textbf{reduced} to G-loc-a & \cite{CascadeHlocClosure}; G-loc-a reduces to Gap G6.2 (day-scale) \\
H-rdr (sufficiency-only) & \textbf{closed} (corollary of P-A) & \cite{CascadeHrdrClosure} \\
\end{tabular}
\end{center}

\textbf{Net.} After the closure work of
\cite{CascadeG64aClosure, CascadeHrdrClosure,
CascadeHlocClosure, CascadeHFICClosure}, the capstone's
hypothesis stack has narrowed substantially: H-meas is a
theorem (G6.4-a closed, 386 cases); H-rdr is closed by
reduction to Theorem P-A; H-loc and H-FIC are reduced to
smaller residual sub-gaps (G-loc-a, G-FIC-a, G-FIC-b), each
lemma-grade. The remaining capstone-introduced hypothesis is
H-lift-fin (sheaf-theoretic bridge refinement), plus the
three lemma-grade residuals. The programme-inherited
hypothesis H-grad is \textbf{frame-level closed at the full
$\widehat{\Z[\varphi]^6}$ level} as of 2026-04-24 via
H-grad-1's sim-closure (builds B1--B5 exact sims, B6
proof-sketch at $G_p$ level; see \cite{CascadeHGrad1}). The
two residuals against full unconditional status are the
cascade-specific $G_2$/octonion pullback $\kappa$ (outlined
but not fully constructed) and the formal equivalence
theorem that $G_p$-canonicality suffices for the full P-A
statement. H-grad remains unconditional at the Fano
$(\Z/2)^3$ level sufficient for all measurement-theoretic
physics (ARIA, Observer rung, Paper~XXXIII). Note that
Theorems~\ref{thm:closure-saturation} and \ref{thm:selection}
inherit full-$\widehat{\Z[\varphi]^6}$-level conditionality
from P-A through the closure-data saturation argument; they
are \emph{not} unconditional at the Fano level alone.

%=====================================================================
\section{The Category of Rung-Structures}
\label{sec:cat-C}
%=====================================================================

We work with a category whose objects are rung-subsets of the
VFD cascade FAN equipped with their inherited closure conditions,
and whose morphisms are closure-preserving inclusions. The
definitions below specialise the standing FAN data of
\cite[\S 1.2]{CascadeObserverQueryAlgebra}.

\subsection{The FAN and its rung set}

\begin{definition}[Cascade FAN, recap from \cite{CascadeCompletenessAudit}]
\label{def:rungs}
Let $\Rungs := \{E_8, H_4, D_4, O, F_{16}, L_{40}, U_0\}$ with
partial order $r \leq r'$ iff $r$ embeds as a Lie-algebra or
Coxeter-group substructure of $r'$. Under this order, $E_8$ is
the top element and the other six rungs are pairwise
incomparable leaves:
\[
    \begin{tikzcd}
    & & E_8 \arrow[ddllll, no head] \arrow[ddll, no head]
        \arrow[dd, no head] \arrow[ddrr, no head]
        \arrow[ddrrrr, no head] \arrow[ddrrrrrr, no head] & & & & \\
    & & & & & & \\
    H_4 & & D_4 & & O & & F_{16} \hspace{1em} L_{40} \hspace{1em} U_0
    \end{tikzcd}
\]
(schematic; the fan has six incomparable leaves below $E_8$,
laid out above with the last three compressed for space).
\end{definition}

\begin{definition}[Closure conditions]
\label{def:closure-cond}
For each $r \in \Rungs$, its closure condition $C_r$ is the
algebraic or geometric constraint characterising rung-$r$ substrate
patches, imported from
\cite[Table in \S 1.3]{CascadeObserverQueryAlgebra}:
\begin{itemize}
\item $C_{E_8}$: icosian construction $(x, \sigma(x)) \in I \times I
  \subset \R^4 \oplus \R^4$;
\item $C_{H_4}$: $H_4$-root action preserves the patch;
\item $C_{D_4}$: $D_4$-Weyl action preserves the patch;
\item $C_{O}$: observer constraint, $\hat o \in S^7$ with octonion
  alternative law (\cite{CascadeObserverQueryAlgebra} E6);
\item $C_{F_{16}}$: $16$-dimensional $\mathrm{Cl}(1,3)$ signature
  selection;
\item $C_{L_{40}}$: icos-$40$ configuration (biological rung);
\item $C_{U_0}$: trivial (all patches collapse to a point).
\end{itemize}
Each $C_r$ determines a subgroupoid $\Gcal_r \subseteq \Gcal$ of
the meta-layer matching groupoid
\cite{CascadeMetaLayer}.
\end{definition}

\subsection{Rung-structures as objects}

\begin{definition}[Rung-structure]
\label{def:rung-structure}
A \emph{rung-structure} is a pair $S = (\underline{S},
\Phi_S)$ where:
\begin{itemize}
\item $\underline{S} \subseteq \Rungs$ is a non-empty up-closed
  subset: if $r \in \underline{S}$ and $r \leq r'$, then $r'
  \in \underline{S}$;
\item $\Phi_S$ is the \emph{installed closure data}: for each
  $r \in \underline{S}$, a specified subgroupoid
  $\Gcal_r^S \subseteq \Gcal_r$ (the rung-$r$ subgroupoid
  accessed by $S$), compatible in the sense that for $r \leq r'$
  with both in $\underline{S}$, the corresponding restriction maps
  $\Gcal_{r'}^S \to \Gcal_r^S$ commute with the FAN's restriction
  structure.
\end{itemize}
We abuse notation and write $S$ for $(\underline{S}, \Phi_S)$
when context is clear.
\end{definition}

\begin{remark}[Up-closure is automatic on the FAN]
\label{rmk:up-closed}
Under the order of Definition~\ref{def:rungs} ($E_8$ top, six
incomparable leaves below), up-closure of
$\underline{S} \subseteq \Rungs$ means: $r \in \underline{S}$
and $r \leq r'$ imply $r' \in \underline{S}$. Because every
leaf is $\leq E_8$, any non-empty $\underline{S}$ containing
at least one leaf automatically contains $E_8$. The up-closure
requirement therefore reduces to: \emph{non-empty
$\underline{S}$ is either $\{E_8\}$, $\{E_8\} \cup L$ for some
non-empty $L \subseteq \text{leaves}$, or $\Rungs$ itself}.
Def.~\ref{def:rung-structure}'s non-emptiness clause excludes
the empty set; no other rung-subset is up-closed.
\end{remark}

\begin{remark}[Role of $\Phi_S$]
\label{rmk:phi-data}
The installed closure data $\Phi_S$ carries \emph{which parts}
of each rung's closure are actually accessed by the rung-
structure $S$. For the maximum-access choice
$\Gcal_r^S = \Gcal_r$ for all $r$, $S$ is the full rung-
structure on $\underline{S}$; restricted choices correspond to
observer instantiations that access only some rung-$r$ patches.
\end{remark}

\subsection{Morphisms}

\begin{definition}[Closure-preserving inclusion]
\label{def:morphism}
A morphism $f \colon S \to S'$ of rung-structures is an inclusion
$\underline{S} \subseteq \underline{S'}$ together with, for each
$r \in \underline{S}$, an inclusion $\Gcal_r^S \hookrightarrow
\Gcal_r^{S'}$ compatible with the FAN restriction maps. The
composition of morphisms is the obvious composition of
inclusions; identities are identity inclusions.
\end{definition}

\begin{definition}[The category $\catC$]
\label{def:catC}
Let $\catC$ denote the category whose objects are rung-structures
(Definition~\ref{def:rung-structure}) and whose morphisms are
closure-preserving inclusions (Definition~\ref{def:morphism}).
\end{definition}

\subsection{Categorical structure of $\catC$}

The following structural properties will be needed to invoke
Ad\'amek's theorem in Section~\ref{sec:existence}.

\begin{proposition}[Terminal object]
\label{prop:terminal}
The category $\catC$ has a terminal object $\mathbf{1}_\catC :=
(\Rungs, \Phi_{\max})$, where $\Phi_{\max}$ assigns $\Gcal_r^{\mathbf 1} =
\Gcal_r$ for each $r$. Equivalently, $\mathbf{1}_\catC$ is the
fully-instantiated rung-structure with full access to every rung.
\end{proposition}

\begin{proof}[Proof sketch]
For any rung-structure $S$, the inclusion $\underline{S}
\subseteq \Rungs$ together with the inclusions $\Gcal_r^S
\hookrightarrow \Gcal_r = \Gcal_r^{\mathbf{1}}$ defines a
morphism $S \to \mathbf{1}_\catC$. Uniqueness of this morphism
follows from the fact that there is only one such inclusion of
each component. Functoriality of the FAN restriction maps makes
the composition compatible.
\end{proof}

\begin{proposition}[Pullbacks along monomorphisms]
\label{prop:pullbacks}
For any $f \colon S \to T$ and $g \colon S' \to T$ in $\catC$
where $g$ is a monomorphism (inclusion as defined in
Definition~\ref{def:morphism}), the pullback $S \times_T S'$
exists and is computed pointwise: $(S \times_T S')_{\mathrm{rung}} =
\underline{S} \cap \underline{S'}$, and
$\Gcal_r^{S \times_T S'} = \Gcal_r^S \cap \Gcal_r^{S'}$ for
each $r \in \underline{S} \cap \underline{S'}$.
\end{proposition}

\begin{proof}[Proof sketch]
The intersection $\Gcal_r^S \cap \Gcal_r^{S'}$ is a subgroupoid
of $\Gcal_r$ because the intersection of subgroupoids is a
subgroupoid. The up-closure of $\underline{S} \cap
\underline{S'}$ is inherited from the up-closure of
$\underline{S}$ and $\underline{S'}$ (intersection of
up-closed sets is up-closed). The universal property of
pullbacks is the usual one for concrete intersection-like
constructions.
\end{proof}

\begin{proposition}[Binary products in $\catC$]
\label{prop:products}
$\catC$ has binary products: for any $S, S' \in \catC$, the
product $S \times S'$ exists and is given by the pullback along
the terminal projections,
\[
    S \times S' \;=\; S \times_{\mathbf{1}_\catC} S',
\]
computed pointwise as
$(S \times S')_{\mathrm{rung}} = \underline{S} \cap
\underline{S'}$ and $\Gcal_r^{S \times S'} = \Gcal_r^S \cap
\Gcal_r^{S'}$ for each $r \in \underline{S} \cap
\underline{S'}$.
\end{proposition}

\begin{proof}
The terminal maps $S \to \mathbf{1}_\catC$ and $S' \to
\mathbf{1}_\catC$ are both monomorphisms
(inclusions into the fully-instantiated rung-structure), so
Proposition~\ref{prop:pullbacks} applies. The pullback along
monomorphisms-to-terminal is, in any category with terminal
object and pullbacks-along-monos, the binary product. The
universal property is inherited from the pullback.
\end{proof}

\begin{corollary}[Finite products and small products of
  $F$-iterates]
\label{cor:finite-products}
Finite products in $\catC$ exist (by iterating
Proposition~\ref{prop:products}), and the countable limit of the
descending $F$-iterate chain $X \times F(X) \times F^2(X) \times
\cdots$ from any $X \in \catC$ exists in $\catC$ by
Proposition~\ref{prop:adamek-sufficient}.
\end{corollary}

\begin{proof}
Immediate from Propositions~\ref{prop:products} and
\ref{prop:adamek-sufficient}.
\end{proof}

\begin{proposition}[Filtered colimits along chains of inclusions]
\label{prop:filtered-colim}
The category $\catC$ admits colimits of $\omega$-chains of
monomorphisms: if $S_0 \hookrightarrow S_1 \hookrightarrow S_2
\hookrightarrow \cdots$ is a chain of inclusions, then the
colimit $S_\infty := \mathrm{colim}_n S_n$ exists with
$\underline{S_\infty} = \bigcup_n \underline{S_n}$ and
$\Gcal_r^{S_\infty} = \bigcup_{n : r \in \underline{S_n}}
\Gcal_r^{S_n}$ for each $r \in \underline{S_\infty}$. Because
$\Rungs$ is finite, the chain stabilises after at most $|\Rungs|
= 7$ strict inclusions.
\end{proposition}

\begin{proof}[Proof sketch]
Union of up-closed sets is up-closed; union of nested
subgroupoids is a subgroupoid (groupoid axioms are closed under
directed unions). Universal property of the colimit is the
standard one. Stabilisation is by finiteness of $\Rungs$.
\end{proof}

\begin{remark}[$\catC$ is essentially finite in the rung-index
  direction, potentially large in the $\Gcal_r$ direction]
\label{rmk:size}
The rung-index $\underline{S}$ ranges over a finite poset of
$\leq 2^{|\Rungs|} = 128$ up-closed subsets. The closure-data
component $\Phi_S$ ranges over subgroupoids of each $\Gcal_r$,
which is a large class. This dual structure is what makes the
self-inquiry functor $F$ (Section~\ref{sec:F}) have non-trivial
dynamics: $F$ acts non-trivially on the $\Gcal_r$ component even
within a fixed $\underline{S}$.
\end{remark}

\subsection{Sufficiency for Ad\'amek}

\begin{proposition}[$\catC$ supports Ad\'amek's theorem]
\label{prop:adamek-sufficient}
$\catC$ has a terminal object (Proposition~\ref{prop:terminal})
and admits $\omega^{\mathrm{op}}$-limits (dual of
Proposition~\ref{prop:filtered-colim} under the rung-finiteness
remark, since descending chains also stabilise). These are the
hypotheses of Ad\'amek's theorem
\cite[Thm.~2.1]{Adamek1974, Rutten2000} for the existence of
a final coalgebra of any endofunctor $F \colon \catC \to
\catC$ preserving $\omega^{\mathrm{op}}$-limits.
\end{proposition}

\begin{proof}[Proof sketch]
Dualise Proposition~\ref{prop:filtered-colim}: a descending
chain $S_0 \hookleftarrow S_1 \hookleftarrow S_2 \hookleftarrow
\cdots$ of inclusions stabilises in the rung-index by finiteness
of $\Rungs$; the $\Gcal_r^{S_n}$ components descend as a nested
intersection $\bigcap_n \Gcal_r^{S_n}$, which is a subgroupoid
(intersection of subgroupoids). The limit is this pointwise
intersection. Universal property is standard. The terminal
object serves as the starting point of the iteration chain.
\end{proof}

%=====================================================================
\section{The Self-Inquiry Functor $F$}
\label{sec:F}
%=====================================================================

This section defines the endofunctor $F : \catC \to \catC$ whose
final coalgebra will be identified with the VFD cascade
(Section~\ref{sec:identification}). Informally, $F(S)$ is the
smallest rung-structure capable of \emph{asking the question
``what is $S$?''}: it must host $S$ as a substructure, contain the
measurement-capable Observer rung $O$, and supply enough closure
data to refer to $S$'s observables as query-algebra elements.

The precise content of the word ``ask'' is supplied by the query
algebra $\Qalg_S \subseteq \Gamma(\Gcal, \Fcal)$ of
\cite[\S\S2--3]{CascadeObserverQueryAlgebra}: $\Qalg_S$ is the
subalgebra of substrate-catalogue sections whose evaluation
requires only the closures $\{C_r : r \in \underline{S}\}$, and
an observer instantiating the rung-set $S$ can measure precisely
those patches whose sheaf sections lie in $\Qalg_S$. In this
language, $F(S)$ is the minimum rung-structure $T$ satisfying
$\Qalg_S \subseteq \Qalg_T$ with $\Qalg_T$ additionally containing
\emph{measurement outcomes} for the $\Qalg_S$-observables — i.e.
$T$ can not only refer to $S$'s observables but resolve them.

\subsection{Why the Observer rung is forced}
\label{subsec:O-forced}

The following is now a theorem of the query-algebra programme
(Gap G6.4 closed), inherited from the Observer-rung measurement-
bridge note:

\begin{proposition}[Observer-rung necessity; $\Qalg_O \cong
  \mathrm{Meas}(S^7, \sigma)$ at theorem-grade after G6.4-a
  closure]
\label{thm:O-necessary}
Import the bridge of \cite{CascadeQOBridge} at theorem-grade,
using the G6.4-a closure of \cite{CascadeG64aClosure} (386
computational Fano-triad sign cases, all pass). For a
rung-structure $T \in \catC$, $\Qalg_T$ contains
measurement-outcome sections (in the sense of
\cite{PaperXXXIII}) iff $O \in \underline{T}$. The Observer-
rung query algebra is isomorphic to the measurement-outcome
algebra, $\Qalg_O \cong \mathrm{Meas}(S^7, \sigma)$, as
$\R$-algebras with involution, via Fano-plane octonion
composition matched to $\sigma$-projection parenthesisation.
\end{proposition}

\begin{proof}[Proof sketch]
The isomorphism $\Qalg_O \cong \mathrm{Meas}(S^7, \sigma)$ is
stated at lemma-grade in \cite{CascadeQOBridge} (working note;
closing G6.4 in the sense that this note records the
isomorphism, with a remaining sign-check sub-gap G6.4-a).
Taking the isomorphism as a standing input, the forward
direction (only-if) follows: if $O \in \underline{T}$, then
$\Qalg_O \subseteq \Qalg_T$
(Proposition~\ref{prop:Q-monotone}), hence $\Qalg_T$ contains
the measurement-outcome sub-algebra.

The converse (if $\Qalg_T$ contains measurement outcomes then
$O \in \underline{T}$): measurement requires the octonion
alternative-law pairing $\Qalg_O \times \Fcal \to \Qalg_O$,
which by the classification in \cite{CascadeQOBridge} is
provided only by the closure condition $C_O$. No other rung's
closure produces this pairing, so the rung-index must contain
$O$. With G6.4-a closed (\cite{CascadeG64aClosure}), this
argument imports \cite{CascadeQOBridge}'s classification at
theorem grade.
\end{proof}

Consequently, any rung-structure from which $S$ can be
\emph{queried}, not merely contained, must satisfy $O \in
\underline{T}$. The self-inquiry functor $F$ enforces this
minimum.

\subsection{Definition of $F$ on objects}
\label{subsec:F-objects}

\begin{definition}[Self-inquiry functor, objects]
\label{def:F-objects}
Let $S = (\underline{S}, \Phi_S) \in \catC$. Its \emph{self-inquiry
closure} is the rung-structure
\[
    F(S) = (\underline{F(S)}, \Phi_{F(S)})
\]
defined by:
\begin{enumerate}
\item \textbf{Underlying rung-set.} The up-closure of
  $\underline{S} \cup \{O\}$ in $\Rungs$:
  \[
      \underline{F(S)} := \mathrm{UC}(\underline{S} \cup \{O\})
      = \underline{S} \cup \{O, E_8\}
  \]
  (the second equality uses $O \leq E_8$ and
  Remark~\ref{rmk:up-closed}).

\item \textbf{Installed closure data.} For each $r \in
  \underline{F(S)}$, specify $\Gcal_r^{F(S)} \subseteq \Gcal_r$
  as follows:
  \begin{enumerate}
  \item For $r \in \underline{S}$ with $r \ne O, E_8$:
    $\Gcal_r^{F(S)} := \Gcal_r^S$ (inherit what $S$ itself
    accesses).
  \item For $r = O$:
    $\Gcal_O^{F(S)} := \Gcal_O^{(S)}$, where $\Gcal_O^{(S)}$ is
    the minimum subgroupoid of $\Gcal_O$ such that the induced
    $\Qalg_O^{F(S)}$ contains a measurement-outcome section for
    every generating observable of $\Qalg_S$
    (Definition~\ref{def:min-O}).
  \item For $r = E_8$: $\Gcal_{E_8}^{F(S)}$ is the subgroupoid
    generated by the FAN-restriction preimages of the
    $\{\Gcal_{r'}^{F(S)} : r' \in \underline{F(S)}, r' \ne E_8\}$
    under the sheaf-compatibility conditions of
    Definition~\ref{def:rung-structure}.
  \end{enumerate}
  If $O \in \underline{S}$ already, the clause in (b) replaces
  $\Gcal_O^S$ by its extension $\Gcal_O^{(S)} \supseteq
  \Gcal_O^S$, which may or may not be a strict extension
  (Remark~\ref{rmk:idempotence-base}).
\end{enumerate}
\end{definition}

\begin{definition}[Minimum measurement-supporting subgroupoid
  $\Gcal_O^{(S)}$ — generator-independent form]
\label{def:min-O}
For $S \in \catC$, define
\[
    \Gcal_O^{(S)} \;:=\; \bigvee_{g \in \Qalg_S}
      \Gcal_O^{(g)} \;\subseteq\; \Gcal_O,
\]
the join (generated subgroupoid) of all $g$-regions ranging
over the \emph{full} subalgebra $\Qalg_S$, not a choice of
generating set. Each $g \in \Qalg_S$ is a sheaf section whose
measurement, via $\sigma$-projection at an octonion direction
$\hat o$, lies in $\Qalg_O$ iff $\hat o$ lies in the closure-
compatible region $\mathcal{O}(g) \subseteq S^7$; its lift to
$\Gcal_O$ generates $\Gcal_O^{(g)}$.
\end{definition}

\begin{remark}[Why the join is over all of $\Qalg_S$]
\label{rmk:gen-independence}
The join over all $g \in \Qalg_S$ is well-defined (join of
subgroupoids is a subgroupoid) and by construction
generator-independent. This removes any dependency of
$\Gcal_O^{(S)}$ on a choice of generating set and makes the
well-definedness of $F$ independent of Gap G6.1 of
\cite{CascadeObserverQueryAlgebra}.
\end{remark}

\begin{hypothesis}[H-lift-fin: finite-cover morphism bridge]
\label{hyp:lift}
For every morphism $\gamma \in \Gcal_O$, there exists a
\emph{finite subset} $\{g_1, \ldots, g_k\} \subset \Qalg_O$
such that $\gamma \in \bigvee_{i=1}^k \Gcal_O^{(g_i)}$.
\end{hypothesis}

\begin{remark}[Derivation route for H-lift-fin — sketch only]
\label{rmk:H-lift-status}
An earlier version of this paper claimed H-lift-fin as a
derived theorem from \cite{CascadeQOBridge}; that claim was
premature. The bridge is an $\R$-algebra-with-involution
isomorphism $\Qalg_O \cong \mathrm{Meas}(S^7, \sigma)$; it
does not by itself give a morphism-to-support correspondence
$\gamma \mapsto \hat o_\gamma$, nor does local compactness of
$S^7$ alone supply a finite cover by observables actually in
$\Qalg_O$ (characteristic functions $\mathbf{1}_{U_i}$ need not
lie in the image of the bridge).

The correct derivation route involves \emph{(i)} establishing a
local observable basis of $\Qalg_O$ indexed by $S^7$-patches
(beyond the algebra-level bridge, this needs a refinement of
the bridge to a \emph{sheaf-theoretic} correspondence), and
\emph{(ii)} a morphism-to-support lemma for $\Gcal_O$ that
supplies the finite cover. Both steps are plausible under the
bridge infrastructure of \cite{CascadeQOBridge} (which is
itself marked ``working-note, lemma-grade'') but require
formalisation beyond the current programme state. H-lift-fin
therefore remains a hypothesis; its derivation is flagged as a
follow-up item in Section~\ref{sec:open}.
\end{remark}

\begin{lemma}[$S^7$-coverage at the terminal, under H-lift-fin]
\label{lem:S7-coverage}
Assume Hypothesis~\ref{hyp:lift} (H-lift-fin). Then
$\Gcal_O^{(\mathbf{1}_\catC)} = \Gcal_O$.
\end{lemma}

\begin{proof}
Fix any morphism $\gamma \in \Gcal_O$. By H-lift-fin, there
exist $g_1, \ldots, g_k \in \Qalg_O \subseteq \Qalg_{\mathbf{1}}$
with $\gamma \in \bigvee_i \Gcal_O^{(g_i)}$. Each
$\Gcal_O^{(g_i)}$ is contained in the global join
$\Gcal_O^{(\mathbf{1}_\catC)} = \bigvee_{g \in
\Qalg_{\mathbf{1}}} \Gcal_O^{(g)}$. Hence $\gamma \in
\Gcal_O^{(\mathbf{1}_\catC)}$. Since $\gamma$ was arbitrary,
$\Gcal_O^{(\mathbf{1}_\catC)} \supseteq \Gcal_O$. The reverse
inclusion $\Gcal_O^{(\mathbf{1}_\catC)} \subseteq \Gcal_O$ is
immediate from Definition~\ref{def:min-O}.
\end{proof}

\begin{remark}[$F(\mathbf{1}_\catC)$ and $F$ of leaves]
\label{rmk:idempotence-base}
When $S = \mathbf{1}_\catC$ (the terminal object, full access),
$\Gcal_O^{(S)} = \Gcal_O$ and $F(\mathbf{1}_\catC) =
\mathbf{1}_\catC$. For $S = (\{U_0, E_8\}, \Phi_{\min})$ (pure
Unity rung), $\Qalg_S$ is the constants and $\Gcal_O^{(S)}$ is
the subgroupoid generated by the identity at any reference
$\hat o_0 \in S^7$; $F(S)$ adjoins the Observer rung minimally,
giving a strictly larger rung-structure than $S$. Both boundary
cases match the intuitive ``$F$ adds observation-capability if
and only if it is missing.''
\end{remark}

\subsection{Definition of $F$ on morphisms}
\label{subsec:F-morphisms}

\begin{definition}[Self-inquiry functor, morphisms]
\label{def:F-morphisms}
Let $f \colon S \to S'$ be a morphism in $\catC$, i.e. an inclusion
$\underline{S} \subseteq \underline{S'}$ together with inclusions
$\Gcal_r^S \hookrightarrow \Gcal_r^{S'}$ for $r \in \underline{S}$.
Define $F(f) \colon F(S) \to F(S')$ by:
\begin{itemize}
\item on the rung-index component:
  $\underline{F(S)} = \underline{S} \cup \{O, E_8\} \subseteq
  \underline{S'} \cup \{O, E_8\} = \underline{F(S')}$;
\item on the closure-data component, for each $r \in
  \underline{F(S)}$, an inclusion
  $\Gcal_r^{F(S)} \hookrightarrow \Gcal_r^{F(S')}$ supplied as
  follows:
  \begin{enumerate}
  \item $r \in \underline{S}$, $r \ne O, E_8$: inherit the
    inclusion $\Gcal_r^S \hookrightarrow \Gcal_r^{S'}$ from $f$;
  \item $r = O$: $\Gcal_O^{(S)} \hookrightarrow \Gcal_O^{(S')}$ is
    the inclusion induced by $\Qalg_S \subseteq \Qalg_{S'}$
    (Proposition~\ref{prop:Q-monotone}): each element $g \in
    \Qalg_S$ maps into $\Qalg_{S'}$, and hence the lift
    $\Gcal_O^{(g)}$ maps into $\Gcal_O^{(g)}$ as a subgroupoid of
    the larger join taken over $\Qalg_{S'}$;
  \item $r = E_8$: inclusion is the restriction-compatibility
    extension of the rung-index inclusions in (a)--(b), using
    Definition~\ref{def:rung-structure}'s closure-compatibility.
  \end{enumerate}
\end{itemize}
\end{definition}

\begin{proposition}[Monotonicity of $\Qalg$]
\label{prop:Q-monotone}
If $f \colon S \hookrightarrow S'$ is a morphism in $\catC$, then
$\Qalg_S \subseteq \Qalg_{S'}$ as subalgebras of
$\Gamma(\Gcal, \Fcal)$.
\end{proposition}

\begin{proof}[Proof sketch]
From Definitions 2.1, 3.1 of
\cite{CascadeObserverQueryAlgebra}: $\Qalg_S$ is generated by
$\{\Qalg_r : r \in \underline{S}\}$ with closure data
$\Gcal_r^S$. Under the inclusion $f$, each generator is mapped
to the corresponding generator of $\Qalg_{S'}$ with the inclusion
$\Gcal_r^S \hookrightarrow \Gcal_r^{S'}$. The generated subalgebra
inclusion follows.
\end{proof}

\subsection{$F$ is an endofunctor}
\label{subsec:F-functor}

\begin{theorem}[Functoriality]
\label{thm:F-functor}
Definitions~\ref{def:F-objects} and~\ref{def:F-morphisms} produce
an endofunctor $F \colon \catC \to \catC$.
\end{theorem}

\begin{proof}[Proof sketch]
\textbf{(a) $F(S) \in \catC$.} The underlying rung-set
$\underline{F(S)}$ is up-closed by construction. Each
$\Gcal_r^{F(S)}$ is a subgroupoid of $\Gcal_r$: for $r \ne O,
E_8$ by hypothesis (inherited from $S$); for $r = O$ by the
join-of-subgroupoids-is-subgroupoid property
(Definition~\ref{def:min-O}); for $r = E_8$ by generated-
subgroupoid construction. Closure-compatibility between rungs
is ensured by the explicit $E_8$-restriction clause.

\textbf{(b) $F(f) \in \catC(F(S), F(S'))$.} The rung-index inclusion
is well-typed. The closure-data inclusions are well-typed
componentwise by Definition~\ref{def:F-morphisms}. Closure-
compatibility with the FAN restriction maps follows from the
commutativity of inclusions across restriction.

\textbf{(c) Identity.} $F(\mathrm{id}_S)$ acts as the identity on
$\underline{F(S)}$ and, componentwise, as identity on each
$\Gcal_r^{F(S)}$. Hence $F(\mathrm{id}_S) = \mathrm{id}_{F(S)}$.

\textbf{(d) Composition.} For $f \colon S \to S'$ and $g \colon S'
\to S''$, the composite $F(g) \circ F(f)$ equals $F(g \circ f)$
componentwise: on the rung-index it is the composite inclusion
$\underline{F(S)} \subseteq \underline{F(S'')}$; on each
closure-data component it is the composite of the individual
inclusions, and these agree with the inclusions $F(g \circ f)$
produces by construction of Definition~\ref{def:F-morphisms}.
\end{proof}

\subsection{Monotonicity}
\label{subsec:F-monotone}

\begin{proposition}[Monotonicity of $F$]
\label{prop:F-monotone}
If $S \hookrightarrow S'$ is a monomorphism (inclusion) in $\catC$,
then $F(S) \hookrightarrow F(S')$ is also a monomorphism. More
strongly, the assignment $S \mapsto F(S)$ is order-preserving with
respect to the closure-preserving-inclusion preorder on $\catC$.
\end{proposition}

\begin{proof}
Immediate from Definition~\ref{def:F-morphisms}: the rung-index
inclusion is preserved, and each closure-data component inclusion
is preserved by construction. Composition with
Theorem~\ref{thm:F-functor}(d) gives the order-preserving
property.
\end{proof}

\begin{corollary}[Idempotence on $\mathbf{1}_\catC$, under
  H-lift-fin]
\label{cor:terminal-fixed}
Assume Hypothesis~\ref{hyp:lift} (H-lift-fin). Then
$F(\mathbf{1}_\catC) = \mathbf{1}_\catC$, so the terminal
object is a (strict) fixed point of $F$.
\end{corollary}

\begin{proof}
The rung-index of $F(\mathbf{1}_\catC)$ adds $\{O, E_8\}$ to
$\underline{\mathbf{1}_\catC} = \Rungs$, leaving it unchanged.
For the closure data: on non-Observer rungs, $F$ inherits
$\Gcal_r^{\mathbf{1}} = \Gcal_r$ directly
(Definition~\ref{def:F-objects}, clause (a)). On the Observer
rung, Lemma~\ref{lem:S7-coverage} (under H-lift-fin) gives
$\Gcal_O^{(\mathbf{1})} = \Gcal_O$. On $E_8$, the FAN-
restriction construction produces $\Gcal_{E_8}^{F(\mathbf{1})}
= \Gcal_{E_8}$ by pulling back the maximal components below.
Hence $F(\mathbf{1}_\catC) = \mathbf{1}_\catC$ componentwise.
\end{proof}

\subsection{Preservation of $\omega^{\mathrm{op}}$-limits}
\label{subsec:F-omega-limits}

This is the critical technical input for the Ad\'amek-theorem
application in Section~\ref{sec:existence}. The one delicate
step is commutation of the generated-subgroupoid operation on
the Observer rung with intersection along descending chains;
we formalise this as a named hypothesis imported from the
query-algebra programme and prove the theorem conditional on
it.

\begin{hypothesis}[H-FIC: Filtered-Intersection-Commutation on
  the Observer and $E_8$ rungs]
\label{hyp:FIC}
For any descending chain $(S_n)_{n \in \N}$ of rung-structures
in $\catC$ with pointwise-intersection limit $S_\infty$, both
of the following hold as subgroupoid identities:
\begin{enumerate}
\item[(H-FIC-O)]
$\Gcal_O^{(S_\infty)} = \bigcap_{n \in \N} \Gcal_O^{(S_n)}$
in $\Gcal_O$.
\item[(H-FIC-$E_8$)]
$\Gcal_{E_8}^{F(S_\infty)} = \bigcap_{n \in \N}
\Gcal_{E_8}^{F(S_n)}$
in $\Gcal_{E_8}$.
\end{enumerate}
\end{hypothesis}

\begin{remark}[H-FIC reduced to G-FIC-a and G-FIC-b via
  support-projection]
\label{rmk:FIC-derivation-route}
\cite{CascadeHFICClosure} reduces H-FIC to smaller residual
sub-gaps via a support-projection construction
$\pi_{S^7} \colon \Gcal_O \to 2^{S^7}$ induced by
Definition~\ref{def:min-O}: each $g \in \Qalg_S$ lifts to
$\Gcal_O^{(g)}$ via its support region $\mathcal{O}(g)
\subseteq S^7$, so $\Gcal_O^{(S)} =
\pi_{S^7}^{-1}(\mathrm{Pow}(\mathcal{O}(\Qalg_S)))$. Preimages
commute with arbitrary intersections (set-theoretic), so
H-FIC-O reduces to an $S^7$-support-region descending-intersection
identity (\textbf{G-FIC-a}), and the rung-finiteness of
$\Rungs$ stabilises the rung-index chain, leaving per-rung
closure-data chains reducible via standard Čech-sheaf
limit-preservation (Kellendonk-Lenz-Savinien 2006). H-FIC-$E_8$
similarly reduces to a generated-subgroupoid-commutes-with-
intersection statement for FAN-restriction preimages
(\textbf{G-FIC-b}, reducible to MacLane Ch.~XII finite-word
factorisation stability). Both G-FIC-a and G-FIC-b are
lemma-grade; estimated closure of both is $1$--$2$ weeks.
\end{remark}

\begin{theorem}[$F$ preserves $\omega^{\mathrm{op}}$-limits,
  under G-FIC-a + G-FIC-b]
\label{thm:F-limits}
Assume residual sub-gaps G-FIC-a and G-FIC-b of
\cite{CascadeHFICClosure} (the reductions of
Hypothesis~\ref{hyp:FIC} via the support-projection argument;
see Remark~\ref{rmk:FIC-derivation-route}). For any descending
chain of monomorphisms
\[
    S_0 \hookleftarrow{} S_1 \hookleftarrow{} S_2
      \hookleftarrow{} \cdots
\]
in $\catC$ with limit $S_\infty := \lim_n S_n$, we have
\[
    F(S_\infty) \;\cong\; \lim_n F(S_n)
\]
canonically in $\catC$.
\end{theorem}

\begin{proof}[Proof]
Work componentwise.

\textbf{(i) Rung-index stabilisation.} Since $\Rungs$ is finite
($|\Rungs| = 7$), the rung-index sequence
$\underline{S_0} \supseteq \underline{S_1} \supseteq \cdots$ is
eventually constant: there exists $N \in \N$ with
$\underline{S_n} = \underline{S_N}$ for all $n \geq N$. Call the
stable value $\underline{S_\infty}$. Then
\begin{align*}
    \underline{F(S_\infty)} &= \underline{S_\infty} \cup \{O, E_8\}
      = \underline{S_N} \cup \{O, E_8\} \\
    &= \underline{F(S_N)} = \lim_n \underline{F(S_n)}
\end{align*}
as eventually-constant sequences.

\textbf{(ii) Closure-data preservation off the Observer rung.} For
$r \in \underline{S_\infty}$ with $r \ne O, E_8$, the closure-data
component is
\[
    \Gcal_r^{F(S_\infty)} = \Gcal_r^{S_\infty}
      = \bigcap_n \Gcal_r^{S_n} = \bigcap_n \Gcal_r^{F(S_n)}
\]
using the defining identity $\Gcal_r^{F(T)} = \Gcal_r^T$ for $r
\ne O, E_8$ (Definition~\ref{def:F-objects}, clause (a)) and
Proposition~\ref{prop:adamek-sufficient} (intersection of
subgroupoids along a descending chain).

\textbf{(iii) Closure-data preservation on the Observer rung.}
For $r = O$, we must show
\[
    \Gcal_O^{(S_\infty)}
    \;=\; \bigcap_n \Gcal_O^{(S_n)}.
\]
By the support-projection reduction of
Remark~\ref{rmk:FIC-derivation-route}, this identity follows
from G-FIC-a (our assumed residual sub-gap): preimages
$\pi_{S^7}^{-1}$ commute with arbitrary intersections, and
G-FIC-a gives the $S^7$-support-region descending-intersection
identity $\mathcal{O}(\Qalg_{S_\infty}) = \bigcap_n
\mathcal{O}(\Qalg_{S_n})$. Combining these yields
$\Gcal_O^{(S_\infty)} = \bigcap_n \Gcal_O^{(S_n)}$.

\textbf{(iv) Closure-data preservation on $E_8$.} The $r = E_8$
case follows from G-FIC-b (our assumed residual sub-gap):
FAN-restriction preimages from leaf rungs commute with
intersection by finite-word-factorisation stability
(\cite{CascadeHFICClosure}), giving $\Gcal_{E_8}^{F(S_\infty)} =
\bigcap_n \Gcal_{E_8}^{F(S_n)}$.

\textbf{(v) Universal property.} The data assembled in (i)--(iv)
is, componentwise, the limit of the diagram $(F(S_n))_n$ in the
category of subgroupoid chains, and its universal property is
inherited pointwise. Any cone over $(F(S_n))_n$ factors uniquely
through the limit cone on each component, hence globally.
\end{proof}

\begin{corollary}[$F$ satisfies the Ad\'amek hypotheses, under
  G-FIC-a + G-FIC-b]
\label{cor:F-adamek-ready}
Assume residual sub-gaps G-FIC-a and G-FIC-b
(\cite{CascadeHFICClosure}). Then $F$ is an endofunctor of
$\catC$ (Theorem~\ref{thm:F-functor}) preserving
$\omega^{\mathrm{op}}$-limits
(Theorem~\ref{thm:F-limits}) on the category $\catC$ that has
$\omega^{\mathrm{op}}$-limits and a terminal object
(Proposition~\ref{prop:adamek-sufficient}). Hence the hypotheses
of Ad\'amek's final-coalgebra theorem
\cite[Thm.~2.1]{Adamek1974}, \cite[Thm.~7.1]{Rutten2000} are
satisfied.
\end{corollary}

\subsection{Examples}
\label{subsec:F-examples}

The following concrete computations ground the definition and
support the identification theorem of
Section~\ref{sec:identification}.

\begin{enumerate}
\item \textbf{$F(\mathbf{1}_\catC) = \mathbf{1}_\catC$}
  (Corollary~\ref{cor:terminal-fixed}). The fully-instantiated
  rung-structure is a strict fixed point; this is the endpoint of
  $F$-iteration starting from any $S \subseteq \mathbf{1}_\catC$.

\item \textbf{$F(\{U_0, E_8\}, \Phi_{\min})$.} Starting from the
  Unity-only rung-structure (trivial closure, constants only),
  $F$ adjoins $O$: $\underline{F(S)} = \{U_0, O, E_8\}$, and
  $\Gcal_O^{(S)}$ is the reference-observer subgroupoid (the
  minimum $\mathcal{O}$-region covering constants).

\item \textbf{$F(\{H_4, E_8\}, \Phi_{QM})$.} A quantum-only
  rung-structure. $F$ adjoins $O$ with $\Gcal_O^{(S)}$ the
  subgroupoid supporting $\sigma$-projection of $H_4$-patches.
  The result is a quantum-measurement-capable structure: the
  minimum configuration for QM measurement, matching the
  two-rung observer case $S = \{O, H_4\}$ of
  \cite[\S3.2, \S6.2]{CascadeObserverQueryAlgebra}.

\item \textbf{$F(\{O, H_4, E_8\}) = \{O, H_4, E_8\}$ strictly
  iff $\Gcal_O^S$ already supports $H_4$-measurement.} When $S$
  already contains $O$, $F(S)$ extends $\Gcal_O^S$ to
  $\Gcal_O^{(S)}$; this is a strict extension iff $\Gcal_O^S$
  does not already cover the $\mathcal{O}$-regions of all
  sections in $\Qalg_S$. The fixed points of $F$ are exactly
  the rung-structures whose $O$-access is already sufficient
  for their own self-querying.
\end{enumerate}

\begin{remark}[The cascade as the grand fixed point]
\label{rmk:cascade-fixed-point}
Section~\ref{sec:identification} will prove that, among all
rung-structures $S$ satisfying $F(S) \cong S$, the canonical
limit of the Ad\'amek chain
$\mathbf{1}_\catC \xleftarrow{} F(\mathbf{1}_\catC)
\xleftarrow{} F^2(\mathbf{1}_\catC) \xleftarrow{} \cdots$
is the $7$-rung VFD cascade FAN of
\cite{CascadeCompletenessAudit}. Example~(4) above shows that
the cascade's self-consistency is a statement about its
$\Gcal_O$-component covering its entire $\Qalg$; this is
precisely the content of the Observer rung as the ``rung that
observes all the rungs below it.''
\end{remark}

\subsection{Discussion}
\label{subsec:F-discussion}

\paragraph{Philosophical reading.} $F$ encodes the act of
self-inquiry: $F(S)$ is the structural minimum for $S$ to be
asked about. Starting from any $S$, iterating $F$ builds the
structural overhead of observing $S$, then of observing the
observer of $S$, and so on; the fixed point is the structure
that observes itself.

\paragraph{Technical scope.} $F$ is defined for all $S \in \catC$
including those not containing $O$; the closure-data clause for
$r = O$ supplies the minimum subgroupoid whether or not $S$
already accesses $O$. The functor and its downstream theorems
use the residual sub-gaps G-FIC-a, G-FIC-b (reductions of
Hypothesis~\ref{hyp:FIC}), Hypothesis~\ref{hyp:lift}
(H-lift-fin), and G-loc-a (reduction of
Hypothesis~\ref{hyp:loc}), together with H-grad and H-meas
imported from \cite{CascadeAccessPrinciple} at full
$\widehat{\Z[\varphi]^6}$ level. H-rdr
(Hypothesis~\ref{hyp:rdr}) is closed by
\cite{CascadeHrdrClosure} and no longer appears as a separate
assumption. The full closure programme is catalogued in
Section~\ref{sec:open}.

\paragraph{Relation to $\mathbf{F1}$.} The scalar axiom $r = 1 +
1/r$ is the numerical shadow of the operator $F$: it is the
1-dimensional analogue in which the ``rung-structure'' is a
single real number and ``self-inquiry closure'' is the map
$r \mapsto 1 + 1/r$. The fixed point of this scalar map is
$\varphi$; the fixed point of $F$ is $\Omegacap$, the full
cascade. Section~\ref{sec:connection} will make this analogy
precise as a scalar-restriction projection $\Omegacap \mapsto
\varphi$.

%=====================================================================
\section{Existence of the Final Coalgebra}
\label{sec:existence}
%=====================================================================

This section applies the classical Ad\'amek final-coalgebra
theorem \cite{Adamek1974, Rutten2000} to the pair $(\catC, F)$ of
Sections~\ref{sec:cat-C}--\ref{sec:F}. The conclusion is that $F$
has a final coalgebra $(\Omegacap, c)$ whose carrier is the
$\omega^{\mathrm{op}}$-limit of the $F$-iteration chain from the
terminal object. Because $F$ fixes the terminal object
(Corollary~\ref{cor:terminal-fixed}), the iteration stabilises
at step $0$, and the final coalgebra is $(\mathbf{1}_\catC,
\mathrm{id})$. The substantive content — that this terminal
object \emph{is} the VFD cascade — is deferred to
Section~\ref{sec:identification}.

\subsection{Ad\'amek's theorem}

We recall the statement in the form of
\cite[Thm.~2.1]{Adamek1974}, \cite[Thm.~7.1]{Rutten2000}.

\begin{theorem}[Ad\'amek final-coalgebra existence]
\label{thm:adamek-statement}
Let $\catD$ be a category with a terminal object $\mathbf{1}$ and
$\omega^{\mathrm{op}}$-limits. Let $G \colon \catD \to \catD$ be
an endofunctor preserving $\omega^{\mathrm{op}}$-limits. Then the
$\omega^{\mathrm{op}}$-chain
\[
    \mathbf{1} \xleftarrow{\,!_{G(\mathbf{1})}\,}
    G(\mathbf{1})
    \xleftarrow{\,G(!_{G(\mathbf{1})})\,}
    G^2(\mathbf{1})
    \xleftarrow{\,G^2(!_{G(\mathbf{1})})\,}
    \cdots
\]
(where $!_X \colon X \to \mathbf{1}$ denotes the unique map to
the terminal object) has a limit $\Omega = \lim_n G^n(\mathbf{1})$,
and $\Omega$ carries a canonical coalgebra structure
$c \colon \Omega \to G(\Omega)$, which is an isomorphism. The
pair $(\Omega, c)$ is the final object in the category of
$G$-coalgebras: for every $G$-coalgebra $(Y, y)$ there exists a
unique coalgebra morphism $u \colon (Y, y) \to (\Omega, c)$.
\end{theorem}

\subsection{Applying the theorem to $(\catC, F)$}

We verify the hypotheses.

\begin{proposition}[$(\catC, F)$ satisfies Ad\'amek's
  hypotheses, under G-FIC-a and G-FIC-b]
\label{prop:adamek-applies}
Assume residual sub-gaps G-FIC-a and G-FIC-b
(\cite{CascadeHFICClosure}). The category $\catC$ has a
terminal object (Proposition~\ref{prop:terminal}) and
$\omega^{\mathrm{op}}$-limits
(Proposition~\ref{prop:adamek-sufficient}). The endofunctor
$F \colon \catC \to \catC$ (Theorem~\ref{thm:F-functor})
preserves $\omega^{\mathrm{op}}$-limits under G-FIC-a +
G-FIC-b (Theorem~\ref{thm:F-limits}).
\end{proposition}

\begin{proof}
Immediate from the cited results.
\end{proof}

\subsection{Computing the iteration chain}

\begin{proposition}[Iteration chain is eventually constant at
  the terminal]
\label{prop:chain-trivial}
The $F$-iteration $\omega^{\mathrm{op}}$-chain from
$\mathbf{1}_\catC$
\[
    \mathbf{1}_\catC \xleftarrow{} F(\mathbf{1}_\catC)
    \xleftarrow{} F^2(\mathbf{1}_\catC) \xleftarrow{} \cdots
\]
is \emph{eventually constant starting at step $0$}:
$F^n(\mathbf{1}_\catC) = \mathbf{1}_\catC$ for all $n \geq 0$,
and every structure map in the chain is the identity on
$\mathbf{1}_\catC$.
\end{proposition}

\begin{proof}
By Corollary~\ref{cor:terminal-fixed},
$F(\mathbf{1}_\catC) = \mathbf{1}_\catC$. Induction on $n$ gives
$F^n(\mathbf{1}_\catC) = \mathbf{1}_\catC$. The structure map
$!_{F(\mathbf{1}_\catC)} \colon F(\mathbf{1}_\catC) \to
\mathbf{1}_\catC$ is the unique morphism from $\mathbf{1}_\catC$
to itself, which is the identity.
\end{proof}

\begin{corollary}[The limit]
\label{cor:limit-terminal}
$\lim_n F^n(\mathbf{1}_\catC) \;\cong\; \mathbf{1}_\catC$
canonically.
\end{corollary}

\begin{proof}
The limit of an eventually-constant chain is the stable value,
with universal property inherited from the identity structure
maps.
\end{proof}

\subsection{The final coalgebra}

\begin{theorem}[Existence of $\Omegacap$, under G-FIC-a +
  G-FIC-b + H-lift-fin]
\label{thm:existence}
Assume residual sub-gaps G-FIC-a and G-FIC-b
(\cite{CascadeHFICClosure}) and
Hypothesis~\ref{hyp:lift} (H-lift-fin). The self-inquiry
functor $F \colon \catC \to \catC$ has a final coalgebra
\[
    \Omegacap := \mathbf{1}_\catC,
    \qquad
    c := \mathrm{id}_{\Omegacap} \colon \Omegacap
    \xrightarrow{\;\cong\;} F(\Omegacap).
\]
For every $F$-coalgebra $(Y, y)$ with $y \colon Y \to F(Y)$,
there exists a unique coalgebra morphism
$u \colon (Y, y) \to (\Omegacap, c)$.
\end{theorem}

\begin{proof}
By Theorem~\ref{thm:adamek-statement} and
Proposition~\ref{prop:adamek-applies}, $F$ has a final coalgebra
whose carrier is $\lim_n F^n(\mathbf{1}_\catC)$. By
Corollary~\ref{cor:limit-terminal} this limit is
$\mathbf{1}_\catC$. The canonical coalgebra structure
$c \colon \Omegacap \to F(\Omegacap)$ is then
$\mathrm{id}_{\mathbf{1}_\catC}$ because $F(\mathbf{1}_\catC) =
\mathbf{1}_\catC$ (Corollary~\ref{cor:terminal-fixed}) and the
canonical map is determined by the universal property of the
limit to agree with the identity structure maps of
Proposition~\ref{prop:chain-trivial}.

Uniqueness of the coalgebra morphism $u \colon (Y, y) \to
(\Omegacap, c)$ follows from terminality of $\mathbf{1}_\catC$
in $\catC$: the unique morphism $!_Y \colon Y \to
\mathbf{1}_\catC = \Omegacap$ in $\catC$ is a coalgebra morphism
because the square
\[
    \begin{tikzcd}
    Y \arrow[r, "y"] \arrow[d, "!_Y"'] &
      F(Y) \arrow[d, "F(!_Y)"] \\
    \Omegacap \arrow[r, "c = \mathrm{id}"'] &
      F(\Omegacap) = \Omegacap
    \end{tikzcd}
\]
commutes: the lower edge is $\mathrm{id}$, and $F(!_Y) \circ y =
!_{F(Y)}$ (unique map to terminal), which equals $!_Y$ since
both are unique maps $Y \to \Omegacap$. Uniqueness of any other
such morphism reduces to terminality of $\Omegacap$ in $\catC$.
\end{proof}

\begin{remark}[Finality reduces to terminality-plus-idempotence]
\label{rmk:finality-terminal}
Theorem~\ref{thm:existence} shows that, for endofunctors that
fix the terminal object, finality of coalgebra and terminality
of carrier coincide. The non-trivial content of the selection-
principle argument is therefore \emph{not} in the existence
step but in the identification step
(Section~\ref{sec:identification}), where we argue that the
terminal object $\mathbf{1}_\catC = (\Rungs, \Phi_{\max})$ is
canonically the VFD cascade of
\cite{CascadeCompletenessAudit}, rather than an abstract
universal construction. That identification is what makes the
cascade uniquely forced by the categorical structure.
\end{remark}

\subsection{Discussion: where the content lives}

\paragraph{What the Ad\'amek step contributes.} The theorem
supplies (i) the \emph{existence} of a final $F$-coalgebra
guaranteed by classical category theory, (ii) the \emph{universal
property} (every $F$-coalgebra maps uniquely into $\Omegacap$),
(iii) the \emph{fixed-point characterisation}
$\Omegacap \cong F(\Omegacap)$, which is what gives the
self-embedding/fractality corollary of
Section~\ref{sec:identification}.

\paragraph{What it does not contribute.} It does not tell us that
$\mathbf{1}_\catC$ is the VFD cascade rather than some other
abstract structure. That requires a concrete identification
theorem: the terminal rung-structure $(\Rungs, \Phi_{\max})$
must be shown isomorphic to the FAN-with-full-closure structure
derived from $F3+F4$ forcing in
\cite{CascadeCompletenessAudit}. Section~\ref{sec:identification}
supplies this identification, which is where the physical
selection-principle content sits.

\paragraph{Comparison with non-trivial cases.} For functors
$G$ with $G(\mathbf{1}) \ne \mathbf{1}$ (the typical case in the
coalgebra literature; e.g.\ $G = (-)^A$ giving streams, or
$G = \mathcal{P}$ giving non-well-founded sets), the Ad\'amek
chain is non-trivial and the final coalgebra is a genuinely new
object. The VFD self-inquiry functor is degenerate in this
respect because the physical substrate is already self-complete
at its top rung — the fully-instantiated rung-structure
\emph{already} questions itself. This is a feature: it is what
\emph{self-closure} of the cascade means categorically. The
non-trivial content is not in the chain but in the concrete
identification of the fixed point with the physical cascade.

\paragraph{An alternative formulation.} One can equivalently
phrase the theorem as: the cascade $\mathbf{1}_\catC$ is the
unique $F$-fixed point in $\catC$ that is a terminal object.
Existence of other $F$-fixed points (Example~(4) of
Section~\ref{subsec:F-examples}) does not conflict with
uniqueness of the \emph{terminal} fixed point, which is what
Ad\'amek's theorem selects. Section~\ref{sec:identification}
will re-derive this from a different angle: as the unique
coalgebra admitting every other coalgebra as a sub-object.

%=====================================================================
\section{Identification: $\Omegacap$ is the VFD Cascade}
\label{sec:identification}
%=====================================================================

This section supplies the physical content that
Theorem~\ref{thm:existence} does not: namely, that the terminal
object $\mathbf{1}_\catC = \Omegacap$ obtained by Ad\'amek's
theorem is canonically isomorphic to the $7$-rung VFD cascade
of \cite{CascadeCompletenessAudit}. The identification separates
cleanly into two parts — \emph{skeleton} (rung-index) and
\emph{closure data} — which are forced by different mechanisms
and must be established separately. We make both explicit to
avoid the circularity of stipulating the full cascade as the
terminal object by fiat.

\subsection{Skeleton: the 7-rung FAN}

The rung-index of the cascade is forced by the completeness
axioms $F3$ and $F4$ established independently of this paper.

\begin{theorem}[Skeleton forcing; recorded in
  \cite{CascadeCompletenessAudit}, theorem source
  \texttt{cascade-foundations.md}]
\label{thm:skeleton-forcing}
Under axioms $\mathbf{F1}$ (self-closure), $F3$ (F1-closure of
rung-recursion), and $F4$ (dimensional-budget conservation),
the rung-index of the physical cascade is the $7$-rung FAN
\[
    \Rungs_{\mathrm{cascade}} \;=\;
    \{E_8, H_4, D_4, O, F_{16}, L_{40}, U_0\}.
\]
No other rung-index satisfies $F3$ and $F4$ simultaneously.
\end{theorem}

This is the input from the completeness programme; it fixes the
underlying rung-set but does \emph{not} fix the closure data.

\subsection{Closure-data saturation}

The closure data of the physical cascade is forced to be
maximal by the access-principle infrastructure, not by $F3+F4$.
We state the saturation condition as a theorem conditional on
the access-principle hypotheses H-grad and H-meas imported
from \cite{CascadeAccessPrinciple}, together with two additional
named hypotheses (H-loc and H-rdr) introduced below.

\begin{hypothesis}[H-loc: witness-separation, weakened form]
\label{hyp:loc}
For every rung $r \in \Rungs$, every morphism
$\gamma \in \Gcal_r$, and every subgroupoid
$\Gcal_r^\Phi \subseteq \Gcal_r$ with $\gamma \notin
\Gcal_r^\Phi$, there exists a section $s_{\gamma,\Phi} \in
\Qalg_{(\Rungs_{\mathrm{cascade}}, \Phi_{\max})}$ such that
$s_{\gamma, \Phi} \notin \Qalg_{(\Rungs_{\mathrm{cascade}},
\Phi)}$. In words: every strict subgroupoid deficiency admits
at least one witness section that distinguishes $\Phi$ from
$\Phi_{\max}$.
\end{hypothesis}

\begin{remark}[Status of H-loc: reduced to G-loc-a]
\label{rmk:H-loc-status}
H-loc is \emph{reduced} in \cite{CascadeHlocClosure} to a
smaller residual sub-gap G-loc-a (rung-overlap avoidance),
using M3d of \cite{CascadeMetaLayer} plus Kellendonk pattern-
equivariant Čech theory. G-loc-a in turn reduces to Gap G6.2
(leaf-rung independence) of
\cite{CascadeObserverQueryAlgebra}. The original M3 of
\cite{CascadeMetaLayer} gives minimality, unique ergodicity,
and linear repetitivity but does \emph{not} directly produce
sections with support on a prescribed
morphism. H-loc is the additional content needed to bridge
from rung-$r$ subgroupoid membership to sheaf-section
membership in $\Qalg_r$. Its proof is a non-trivial extension
of M3 left for the query-algebra-presentations follow-up work
(\cite[Gap G6.1]{CascadeObserverQueryAlgebra}).
\end{remark}

\begin{lemma}[Section deficiency from closure-data restriction,
  under H-loc]
\label{lem:char-restriction}
Assume Hypothesis~\ref{hyp:loc} (H-loc). Let $S =
(\Rungs_{\mathrm{cascade}}, \Phi) \in \catC$ be a full-skeleton
rung-structure with closure data $\Phi$. If there is some $r
\in \Rungs_{\mathrm{cascade}}$ with $\Gcal_r^\Phi \subsetneq
\Gcal_r$ (strict deficiency), then there exists a section $s
\in \Gamma(\Gcal, \Fcal)$ with $s \in
\Qalg_{(\Rungs_{\mathrm{cascade}}, \Phi_{\max})}$ but $s
\notin \Qalg_{(\Rungs_{\mathrm{cascade}}, \Phi)}$. Equivalently,
$\Qalg_{(\Rungs_{\mathrm{cascade}}, \Phi)} \subsetneq
\Qalg_{(\Rungs_{\mathrm{cascade}}, \Phi_{\max})}$
(no claim is made about $\Qalg_{(\Rungs_{\mathrm{cascade}},
\Phi_{\max})} = \Gamma(\Gcal, \Fcal)$; that is the content of
H-rdr, used separately).
\end{lemma}

\begin{proof}
By the strict deficiency hypothesis, there exists a morphism
$\gamma \in \Gcal_r \setminus \Gcal_r^\Phi$. By H-loc
(witness-separation form), there is a section
$s_{\gamma,\Phi} \in
\Qalg_{(\Rungs_{\mathrm{cascade}}, \Phi_{\max})}$ with
$s_{\gamma, \Phi} \notin
\Qalg_{(\Rungs_{\mathrm{cascade}}, \Phi)}$. Setting $s :=
s_{\gamma,\Phi}$ proves the lemma.
\end{proof}

\begin{remark}[What this lemma does and does not say]
\label{rmk:char-restriction-scope}
Lemma~\ref{lem:char-restriction} establishes only the weaker
statement \emph{some section is lost} under strict closure-
data restriction. It does not attempt the stronger character-
class-level statement ``some whole $\Fcal(\chi)$ is missing'':
a strict subgroupoid deficiency may remove a proper sub-
vector-space within each character component rather than a
whole character class. The weaker statement suffices for the
closure-saturation theorem below.
\end{remark}

\begin{hypothesis}[H-rdr: sufficiency-direction only]
\label{hyp:rdr}
The Access Principle \cite[Theorem P-A]{CascadeAccessPrinciple}
extends from rung-sets $S \subseteq \Rungs$ to rung-data-
refined objects $(S, \Phi)$ in the following
\emph{sufficiency} sense: for $\Phi = \Phi_{\max}$ (the full-
access closure data),
$\Qalg_{(\Rungs_{\mathrm{cascade}}, \Phi_{\max})} =
\Gamma(\Gcal, \Fcal)$. No strict-monotonicity or reverse
implication is claimed in H-rdr; the necessity direction is
established separately using H-loc via
Lemma~\ref{lem:char-restriction}.
\end{hypothesis}

\begin{remark}[Status of H-rdr: closed]
\label{rmk:H-rdr-status}
H-rdr is now \emph{closed} by \cite{CascadeHrdrClosure}: the
refined query algebra
$\Qalg_{(\Rungs_{\mathrm{cascade}}, \Phi_{\max})}$ is
definitionally identical to the bare-rung
$\Qalg_{\Rungs_{\mathrm{cascade}}}$ of
\cite{CascadeAccessPrinciple}, so the sufficiency direction
H-rdr reduces trivially to Theorem P-A at $S =
\Rungs_{\mathrm{cascade}}$ with no new analytic content. We
retain the hypothesis label \texttt{hyp:rdr} for historical
bookkeeping; downstream theorems no longer require H-rdr as
a separate assumption.
\end{remark}

\begin{theorem}[Closure-data saturation; conditional on H-loc
  (reduced to G-loc-a)]
\label{thm:closure-saturation}
Assume Hypothesis~\ref{hyp:loc} (H-loc, witness-separation form;
reduced in \cite{CascadeHlocClosure} to the residual sub-gap
G-loc-a which in turn reduces to Gap G6.2 of
\cite{CascadeObserverQueryAlgebra}). H-rdr is now closed as a
corollary of Theorem P-A (\cite{CascadeHrdrClosure}), so H-grad
and H-meas enter upstream via P-A without separate hypothesis
status. The full-access
rung-structure $(\Rungs_{\mathrm{cascade}}, \Phi_{\max})$ is
the unique rung-structure in $\catC$ with underlying rung-index
$\Rungs_{\mathrm{cascade}}$ whose query algebra coincides with
the full meta-layer section algebra:
$\Qalg_{(\Rungs_{\mathrm{cascade}}, \Phi)} =
\Gamma(\Gcal, \Fcal)$ iff $\Phi = \Phi_{\max}$ (up to closure-
preserving isomorphism).
\end{theorem}

\begin{proof}
\textbf{Sufficiency} ($\Phi = \Phi_{\max} \Rightarrow \Qalg =
\Gamma(\Gcal, \Fcal)$): closed in \cite{CascadeHrdrClosure}
by reduction of H-rdr to Theorem P-A at $S =
\Rungs_{\mathrm{cascade}}$ with full closure data. No
additional hypothesis beyond H-grad and H-meas (imported via
P-A) is needed for this direction.

\textbf{Necessity} ($\Phi \ne \Phi_{\max} \Rightarrow \Qalg
\subsetneq \Gamma(\Gcal, \Fcal)$): if $\Phi \ne \Phi_{\max}$
then there exists $r$ with $\Gcal_r^\Phi \subsetneq \Gcal_r$.
By Lemma~\ref{lem:char-restriction} (under H-loc), there is a
section $s \in \Gamma(\Gcal, \Fcal) \setminus
\Qalg_{(\Rungs_{\mathrm{cascade}}, \Phi)}$. Hence the
inclusion is strict.

\textbf{Uniqueness} (up to closure-preserving isomorphism):
the biconditional established by sufficiency and necessity
determines $\Phi$ (as a closure-data assignment) from the
condition $\Qalg = \Gamma(\Gcal, \Fcal)$ — specifically,
$\Phi$ must equal $\Phi_{\max}$ because necessity rules out
every other $\Phi$.
\end{proof}

\begin{definition}[The VFD cascade]
\label{def:cascade-structure}
The \emph{VFD cascade} is the rung-structure
\[
    \mathrm{Cascade} \;:=\; (\Rungs_{\mathrm{cascade}},
      \Phi_{\max})
\]
where $\Rungs_{\mathrm{cascade}}$ is the rung-index forced by
Theorem~\ref{thm:skeleton-forcing} and $\Phi_{\max}$ is the
closure data forced (under H-grad and H-meas, imported via
P-A, plus H-loc in its witness-separation form, reduced in
\cite{CascadeHlocClosure} to G-loc-a) by
Theorem~\ref{thm:closure-saturation}.
\end{definition}

\begin{remark}[Why the definition is not circular]
\label{rmk:cascade-not-circular}
Definition~\ref{def:cascade-structure} is a \emph{composition}
of two independent forcing theorems: $F3+F4$ force the
skeleton (Theorem~\ref{thm:skeleton-forcing}, unconditional),
and the hypothesis stack (H-grad, H-meas, H-loc, H-rdr) forces
$\Phi_{\max}$ (Theorem~\ref{thm:closure-saturation},
conditional). Neither forcing uses the terminal object of
$\catC$; both are programme-inherited. The subsequent
identification of this object with $\mathbf{1}_\catC$
(Theorem~\ref{thm:selection}) is therefore a non-trivial
structural claim rather than a tautology.
\end{remark}

\subsection{The selection-principle theorem}

\begin{theorem}[Selection principle; conditional on the full
  hypothesis stack]
\label{thm:selection}
Assume residual sub-gaps G-FIC-a and G-FIC-b
(\cite{CascadeHFICClosure}, reductions of H-FIC per
Remark~\ref{rmk:FIC-derivation-route});
Hypothesis~\ref{hyp:lift} (H-lift-fin); and residual sub-gap
G-loc-a (\cite{CascadeHlocClosure}, reduction of H-loc;
further reduces to Gap G6.2). H-rdr is closed by reduction to
Theorem P-A (\cite{CascadeHrdrClosure}) and is no longer a
separate hypothesis. H-grad and H-meas import from
\cite{CascadeAccessPrinciple}; H-meas is theorem-grade via
G6.4-a closure (\cite{CascadeG64aClosure}); H-grad-Fano is
unconditional (\cite{CascadeFanoLift}). \emph{The theorem
inherits full-$\widehat{\Z[\varphi]^6}$-level conditionality
through P-A} — Fano-level H-grad alone does not suffice for
the closure-saturation step of
Theorem~\ref{thm:closure-saturation}; the full H-grad-1
Conway-Sloane identification is required at full P-A. The
remaining capstone-introduced hypothesis is H-lift-fin plus
the lemma-grade residuals G-FIC-a, G-FIC-b, G-loc-a. Under
this stack, there is a canonical isomorphism
\[
    \iota \colon \mathrm{Cascade}
    \xrightarrow{\;\cong\;} \Omegacap
\]
of rung-structures in $\catC$, compatible with the coalgebra
structure: $c \circ \iota = F(\iota) \circ c_{\mathrm{cascade}}$,
where $c_{\mathrm{cascade}} = \mathrm{id}_{\mathrm{Cascade}}$.
Consequently, the VFD cascade is simultaneously the unique
final $F$-coalgebra and the substrate-maximal rung-structure
over the skeleton $\Rungs_{\mathrm{cascade}}$.
\end{theorem}

\begin{proof}
$\mathbf{1}_\catC$ is the terminal object
(Proposition~\ref{prop:terminal}) with rung-index $\Rungs$ and
closure data $\Phi_{\max}$. $\mathrm{Cascade}$ has rung-index
$\Rungs_{\mathrm{cascade}}$ (Theorem~\ref{thm:skeleton-forcing})
and closure data $\Phi_{\max}$
(Theorem~\ref{thm:closure-saturation}, under the stated
hypotheses). The rung-index coincidence
$\Rungs_{\mathrm{cascade}} = \Rungs$ is a \emph{structural}
match between two independently-derived results: $F3+F4$ gives
the specific $7$-rung set, and $\Rungs$ is by definition that
same $7$-rung set in the underlying category. The closure-data
coincidence $\Phi_{\max} = \Phi_{\max}$ holds trivially once
both sides have been forced to be maximal. The canonical
isomorphism $\iota$ is the identity on both components.

Compatibility with the coalgebra structure: by
Theorem~\ref{thm:existence} the coalgebra on $\Omegacap =
\mathbf{1}_\catC$ is $\mathrm{id}_{\Omegacap}$; by
Corollary~\ref{cor:terminal-fixed}, $F(\mathbf{1}_\catC) =
\mathbf{1}_\catC$. Under the structural identification $\iota$
(established above, \emph{before} invoking coalgebra data),
$\mathrm{Cascade}$ is literally $\mathbf{1}_\catC$, hence
$F(\mathrm{Cascade}) = \mathrm{Cascade}$ follows by
transporting the terminal-object identity along $\iota$. The
compatibility square $c \circ \iota = F(\iota) \circ
c_{\mathrm{cascade}}$ then reduces to $\iota = F(\iota)$,
which is the identity. Note the non-circularity: the
identification $\iota$ is given by rung-index and closure-data
coincidence (depends only on Theorems~\ref{thm:skeleton-forcing}
and \ref{thm:closure-saturation}, not on coalgebra data), and
the F-fixedness of $\mathrm{Cascade}$ is then inherited from
that of $\mathbf{1}_\catC$.

The content of the theorem is that two independently-derived
structural descriptions — \emph{(a)} $(F3+F4)$-forced skeleton
combined with access-principle-forced saturation, and
\emph{(b)} Ad\'amek's final coalgebra of $F$ — coincide as the
same rung-structure. This coincidence is where the capstone's
selection-principle content resides.
\end{proof}

\begin{remark}[Why this is the selection principle]
\label{rmk:why-selection}
The classical selection-principle problem is: \emph{why this
structure and not another?} Theorem~\ref{thm:selection} answers
this for the VFD cascade in three equivalent forms:
\begin{enumerate}
\item \textbf{Forcing form.} $F3+F4$ of
  \cite{CascadeCompletenessAudit} force the $7$-rung FAN under
  the axiom $\mathbf{F1}$ and a dimensional-budget argument.
\item \textbf{Coalgebraic form
  (Theorem~\ref{thm:existence}).} $F$ has a unique final
  coalgebra, which is its terminal carrier.
\item \textbf{Identified form
  (Theorem~\ref{thm:selection}).} These two descriptions are
  the same object, so the cascade is simultaneously forced by
  $F3+F4$ and selected by the categorical universal property.
\end{enumerate}
Any of the three forms individually is a partial answer. The
capstone's value is in showing they are a single answer.
\end{remark}

\subsection{Corollaries}

\begin{corollary}[Universal property]
\label{cor:universal-property}
For any rung-structure $Y \in \catC$ equipped with an
$F$-coalgebra structure $y \colon Y \to F(Y)$, there exists a
unique coalgebra morphism $u \colon (Y, y) \to
(\mathrm{Cascade}, \mathrm{id})$. In particular, every
``self-questioning'' rung-structure maps uniquely into the
cascade.
\end{corollary}

\begin{proof}
Theorem~\ref{thm:existence} plus Theorem~\ref{thm:selection}:
the unique coalgebra morphism $(Y, y) \to (\Omegacap, c)$
transported along $\iota$.
\end{proof}

\begin{corollary}[Fixed-point self-embedding]
\label{cor:fixed-point}
$\mathrm{Cascade} \cong F(\mathrm{Cascade})$ canonically. In
particular, the cascade is a fixed point of the self-inquiry
functor up to isomorphism, and in fact up to the identity.
\end{corollary}

\begin{proof}
Corollary~\ref{cor:terminal-fixed} plus
Theorem~\ref{thm:selection}.
\end{proof}

\begin{corollary}[Terminality + $F$-fixedness]
\label{cor:min-max}
$\mathrm{Cascade}$ is simultaneously:
\begin{itemize}
\item \textbf{Terminal in $\catC$}: for every $Y \in \catC$
  there is a unique morphism $Y \to \mathrm{Cascade}$
  (Proposition~\ref{prop:terminal} plus
  Theorem~\ref{thm:selection});
\item \textbf{$F$-fixed}: $F(\mathrm{Cascade}) =
  \mathrm{Cascade}$ (Corollary~\ref{cor:terminal-fixed}).
\end{itemize}
Together these two properties give the final-$F$-coalgebra
universal property of Corollary~\ref{cor:universal-property}.
\end{corollary}

\subsection{Fractality: precise and imprecise readings}

\begin{corollary}[Self-embedding form of fractality]
\label{cor:fractality}
By Corollary~\ref{cor:fixed-point}, $\mathrm{Cascade} \cong
F(\mathrm{Cascade})$; iterating, $\mathrm{Cascade} \cong
F^n(\mathrm{Cascade})$ for every $n \in \N$. Every rung
$r \in \Rungs$ that appears in the cascade also appears in
$F^n(\mathrm{Cascade})$ with the same closure subgroupoid
$\Gcal_r$; the cascade is therefore structurally stable under
iterated self-inquiry.
\end{corollary}

\begin{remark}[What fractality does and does not mean here]
\label{rmk:fractality-caveat}
Corollary~\ref{cor:fractality} gives a categorical self-embedding
statement: the cascade embeds into its self-inquiry closure as
the identity. It does \emph{not} on its own supply the stronger
``each rung contains a scaled copy of the whole cascade'' reading
sometimes associated with the word ``fractal.'' That stronger
statement — which can be expected to hold for appropriate
substrate scales via $\mathbf{F1}$'s golden-ratio self-similarity
— requires an additional rung-wise refinement structure
(scale-dilation, phason-scale recursion) that we do not set up
here. This paper's fractality corollary is the weaker, cleaner,
categorical one; we flag the stronger reading as an open item in
Section~\ref{sec:open}.
\end{remark}

\subsection{Spinozist reading}

\begin{remark}[Mathematical necessitarianism]
\label{rmk:spinoza}
Theorem~\ref{thm:selection} is the mathematical content of
Spinoza's thesis ``everything follows from the nature of things
with absolute necessity'' (\emph{Ethics} I, Prop.~29), applied
to the structural skeleton of physics. The cascade is not a
chosen substrate; it is selected by a categorical universal
property under mild structural hypotheses on rung-structures.
Einstein's question ``did the universe have any choice?'' is
answered in the negative at the level of the
rung-structural skeleton, while remaining silent on the
independent question of why the closure data $\Phi_{\max}$ takes
its specific algebraic forms (the latter is the content of the
supporting papers in the programme). Section~\ref{sec:philosophy}
frames this more carefully and flags what the result does not
claim.
\end{remark}

%=====================================================================
\section{The Four-Phase Comonad $(T, \varepsilon, \delta)$}
\label{sec:comonad}
%=====================================================================

Section~\ref{sec:identification} settled the \emph{structural}
content of the selection principle: the VFD cascade is the
terminal object and final $F$-coalgebra of $\catC$. This section
settles the \emph{dynamical} content: the four VFD operations —
decoherence / closure-memory / measurement-projection /
crystallisation — fit together as a comonad $(T, \varepsilon,
\delta)$ arising from $F$ via the standard cofree-comonad
construction of \cite[\S 3.5]{Jacobs2017}. The identification
converts a previously-scattered collection of dynamical laws
across Papers XXI, XXXI, XXXIII, XLIV, ACT and VFD
Crystallisation into one categorical structure.

\subsection{The cofree comonad on $F$}

\begin{definition}[Cofree comonad, after
  \cite{Jacobs2017}~\S 3.5]
\label{def:cofree-comonad}
For the endofunctor $F \colon \catC \to \catC$ preserving
$\omega^{\mathrm{op}}$-limits (Theorem~\ref{thm:F-limits}), the
\emph{cofree comonad on $F$} is the endofunctor $T \colon
\catC \to \catC$ defined, for each $X \in \catC$, as the final
coalgebra of the functor
\[
    G_X \colon \catC \to \catC, \qquad
    G_X(Z) := X \times F(Z),
\]
with natural transformations $\varepsilon \colon T \Rightarrow
\mathrm{id}_\catC$ and $\delta \colon T \Rightarrow T \circ T$
determined by the universal properties of the finality of
$T(X)$.
\end{definition}

In our setting this construction simplifies because $F$ is
idempotent on the terminal: $F(\mathbf{1}_\catC) =
\mathbf{1}_\catC$ (Corollary~\ref{cor:terminal-fixed}). For
$X = \mathbf{1}_\catC$, $T(\mathbf{1}_\catC)$ is the final
coalgebra of $Z \mapsto \mathbf{1}_\catC \times F(Z)$, which
reduces to the final coalgebra of $Z \mapsto F(Z)$ itself (the
product with a terminal is the identity), and this is
$\mathbf{1}_\catC$ by Theorem~\ref{thm:existence}. Hence
\[
    T(\mathbf{1}_\catC) \;=\; \mathbf{1}_\catC,
    \qquad
    T(\mathrm{Cascade}) \;=\; \mathrm{Cascade}.
\]

\begin{remark}[$T$ on non-terminal objects]
\label{rmk:T-non-terminal}
For $X \in \catC$ strictly smaller than $\mathbf{1}_\catC$,
$T(X)$ is non-trivial: it is the
$\omega^{\mathrm{op}}$-limit of the chain
\[
    X \xleftarrow{\pi_1} X \times F(X) \xleftarrow{\pi_1}
    X \times F(X \times F(X)) \xleftarrow{\pi_1} \cdots
\]
and records all of $X$'s observable iterated closures. The
projections $\pi_1$ in the chain are the iterated
``context-forgetting'' maps; $T(X)$ is the total
context-retention object.
\end{remark}

\subsection{The four VFD operations as comonad components}
\label{subsec:four-phase-id}

We now identify each piece of the comonad with an existing VFD
construction. The identification is a claim about correspondence
at the structural level; the detailed correspondence proofs
cite the originating papers.

\subsubsection*{(a) The functor $T$ as context-retention /
  decoherence}

$T \colon \catC \to \catC$ sends a rung-structure $X$ to its
context-retention object $T(X)$. The projection $\varepsilon_X
\colon T(X) \to X$ is the ``drop all context, keep only $X$''
map; the full $T(X)$ retains $X$ together with all iterated
$F$-closures, which is precisely the data needed to track
coarsened-sector / decoherent-state structure.

\begin{claim}[Forget-phase identification]
\label{claim:T-decoherence}
Under the identification $X \mapsto $ ``quantum sector observed
at rung-structure $X$,'' $T$ corresponds to the
$\sigma$-projection / decoherence coarsening of
\cite{PaperXXXI, PaperXLIV}: $T(X)$ is the space of
context-retained sectors, $\varepsilon_X$ is the ``trace out
environment and keep the system at $X$'' projection.
\end{claim}

\subsubsection*{(b) The comultiplication $\delta$ as
  closure-memory / Remember}

The comultiplication $\delta_X \colon T(X) \to T(T(X))$ encodes
``context of context'': given a context-retained object $T(X)$,
$\delta_X$ lifts it to context-of-context $T^2(X)$, i.e.\ adds
\emph{memory} of the context retention.

\begin{claim}[Remember-phase identification]
\label{claim:delta-memory}
Under the identification of $T$ with context-retention, $\delta$
is \emph{proposed} to correspond to the learned-$W$ memory
upgrade of \cite{ACT-L1} and the closure-dynamical structure of
\cite{PaperXXI}: context-retention $T(X)$ is upgraded to
context-of-context $T^2(X)$ via installation of a learned
weight / memory kernel on the retained context. The specific
mapping is interpretive; the rigorous line-by-line match is
scope for the follow-up paper.
\end{claim}

\subsubsection*{(c) $T^2$ as measurement / Experience}

The twice-iterated context object $T^2(X)$ is the observable
structure that records context-of-context. In VFD operational
terms, this is exactly measurement in the Paper XXXIII sense:
measurement selects a closure-projection of the observable
sheaf section, and the result is an object whose context
includes the context of the measurement itself.

\begin{claim}[Experience-phase identification]
\label{claim:T2-measurement}
$T^2(X)$ corresponds to measurement-as-closure-projection of
\cite{PaperXXXIII}: $T^2(X)$ is the space of
doubly-context-retained sectors, which is the minimum structure
needed to represent ``observer observing system observing
itself.'' The measurement outcome is the value of a
$T^2$-observable, and the Paper XXXIII measurement dynamics
factors through the $T^2$-structure.
\end{claim}

\subsubsection*{(d) The counit $\varepsilon$ as crystallisation /
  Answer}

The counit $\varepsilon_X \colon T(X) \to X$ is the ``crystallise
context to answer'' map: given context-retained object $T(X)$,
$\varepsilon_X$ reduces it to a specific $X$-valued answer,
collapsing retained context to a decision.

\begin{claim}[Answer-phase identification]
\label{claim:epsilon-crystallisation}
Under the identification $X = $ ``closure-functional target,''
$\varepsilon$ corresponds to crystallisation via
$\Fcrystal = \alpha R + \beta E - \gamma Q$ minimisation
\cite{VFDCrystallisation}: $T(X)$'s context is crystallised to
the $\Fcrystal$-minimal $X$-section.
\end{claim}

\subsection{Comonad laws (categorical part, conditional on G-FIC-a + G-FIC-b)}

The comonad identities
\begin{align*}
    \varepsilon_T \circ \delta &= \mathrm{id}_T
      \qquad \text{(left counit)} \\
    T \varepsilon \circ \delta &= \mathrm{id}_T
      \qquad \text{(right counit)} \\
    T \delta \circ \delta &= \delta_T \circ \delta
      \qquad \text{(coassociativity)}
\end{align*}
are automatic consequences of the cofree-comonad construction
(Definition~\ref{def:cofree-comonad}), with products supplied by
Proposition~\ref{prop:products}. We state this categorical
content as a theorem of standard coalgebra, then separate the
VFD-interpretation claim into its own result.

\begin{theorem}[$(T, \varepsilon, \delta)$ is a comonad on
  $\catC$]
\label{thm:comonad-laws}
Under residual sub-gaps G-FIC-a and G-FIC-b
(\cite{CascadeHFICClosure}, so that $F$ preserves
$\omega^{\mathrm{op}}$-limits by Theorem~\ref{thm:F-limits}),
the cofree-comonad construction on $F$ yields a comonad
$(T, \varepsilon, \delta)$ on $\catC$. The three comonad laws
(left counit, right counit, coassociativity) hold as
equalities of natural transformations.
\end{theorem}

\begin{proof}
We need (i) existence of the final coalgebra of $G_X(Z) = X
\times F(Z)$ for each $X \in \catC$, and (ii) the resulting
$(T, \varepsilon, \delta)$ satisfies the comonad laws.

For (i), we verify directly that $G_X = (X \times -) \circ F$
preserves $\omega^{\mathrm{op}}$-limits in our specific $\catC$.
Let $(Y_n)_{n \in \N}$ be a descending chain with
pointwise-intersection limit $Y_\infty$ (existence by
Proposition~\ref{prop:adamek-sufficient}). By
Proposition~\ref{prop:products}, in $\catC$ the binary product
is computed pointwise as intersection:
\[
    (X \times Y)_{\mathrm{rung}} = \underline{X} \cap
    \underline{Y}, \qquad
    \Gcal_r^{X \times Y} = \Gcal_r^X \cap \Gcal_r^Y.
\]
For the chain $(X \times Y_n)_{n \in \N}$, the pointwise
intersections
\[
    \bigcap_n (\underline{X} \cap \underline{Y_n})
    = \underline{X} \cap \bigcap_n \underline{Y_n}
    = \underline{X} \cap \underline{Y_\infty}
\]
and
\[
    \bigcap_n (\Gcal_r^X \cap \Gcal_r^{Y_n})
    = \Gcal_r^X \cap \bigcap_n \Gcal_r^{Y_n}
    = \Gcal_r^X \cap \Gcal_r^{Y_\infty}
\]
use only the basic set-theoretic identity $A \cap \bigcap_n
B_n = \bigcap_n (A \cap B_n)$, componentwise. Hence $\lim_n
(X \times Y_n) = X \times Y_\infty = X \times \lim_n Y_n$; i.e.\
$X \times -$ preserves $\omega^{\mathrm{op}}$-limits in $\catC$.

Composing with $F$, which preserves
$\omega^{\mathrm{op}}$-limits by Theorem~\ref{thm:F-limits},
gives that $G_X$ preserves $\omega^{\mathrm{op}}$-limits.
Ad\'amek's theorem (Theorem~\ref{thm:adamek-statement}) then
yields a final coalgebra of $G_X$; call its carrier $T(X)$. The
assignment $X \mapsto T(X)$ extends to an endofunctor of
$\catC$ by naturality of the final-coalgebra construction in
$X$.

For (ii), by \cite[Prop.~3.5.7, Cor.~3.5.8]{Jacobs2017}, the
cofree comonad on $F$ — which is exactly $T$ as constructed
above — carries natural transformations $\varepsilon \colon T
\Rightarrow \mathrm{id}_\catC$ (the counit) and $\delta \colon
T \Rightarrow T^2$ (the comultiplication) determined by the
finality of $T(X)$, with the three comonad laws holding as
equalities of natural transformations.
\end{proof}

\subsection{VFD-operation identifications (interpretive claim)}

Theorem~\ref{thm:comonad-laws} gives us a comonad $(T,
\varepsilon, \delta)$ on $\catC$ with automatic laws. The
substantive programme claim is that the four VFD operations
realise this comonad:

\begin{claim}[Four-phase realisation of the comonad]
\label{claim:four-phase-realisation}
The natural transformations
$(T, \varepsilon, \delta)$ of Theorem~\ref{thm:comonad-laws}
correspond to the four VFD dynamical operations under the
identifications of Section~\ref{subsec:four-phase-id}:
\begin{center}
\begin{tabular}{l|l}
Categorical component & VFD operation \\ \hline
$T$ (context-retention) & Forget (decoherence $\sigma$-projection) \\
$\delta$ (comultiplication) & Remember (closure-memory lift) \\
$T^2$ (context-of-context) & Experience (measurement-as-closure-projection) \\
$\varepsilon$ (counit) & Answer (crystallisation via
$\Fcrystal = \alpha R + \beta E - \gamma Q$) \\
\end{tabular}
\end{center}
\end{claim}

\begin{remark}[Status of Claim~\ref{claim:four-phase-realisation}]
\label{rmk:realisation-status}
The categorical content (Theorem~\ref{thm:comonad-laws}) is
conditional on G-FIC-a + G-FIC-b only. The identification of
the categorical natural transformations with the four VFD
operations is an \emph{interpretive claim} rather than a
theorem in the present paper: it says that the specific VFD
dynamics of \cite{PaperXXI, PaperXXXI, PaperXXXIII, PaperXLIV,
ACT-L1, VFDCrystallisation} can each be read as the appropriate
component of the comonad.

We do not attempt a line-by-line verification of this
identification here; doing so rigorously requires
expressing each VFD operation in the specific category $\catC$
and checking equality against the automatic comonad-structure
maps. Such a verification is the content of a dedicated follow-
up paper (flagged in Section~\ref{sec:open}, item~4), not of
the capstone. What the capstone establishes is that the
\emph{categorical scaffolding} into which the four phases must
fit exists and has the expected laws; the fit itself is then
a localised check in each supporting paper, not a new
theorem.
\end{remark}

\subsection{The four-phase cycle as a structured dynamics}

\begin{remark}[Reading the comonad as a cycle]
\label{rmk:four-phase-cycle}
The three comonad identities together assert that the sequence
Forget $\to$ Remember $\to$ Experience $\to$ Answer $\to$ Forget
(next iteration) is consistent: the dynamics returns to its
starting context with the memory of the cycle accumulated via
$\delta$. The counit closes the cycle (Answer returns to the
starting context-retained state); the coassociativity ensures
the memory upgrade commutes with itself, so multi-step cycles
compose coherently.

In operational terms: the VFD substrate continuously (i) retains
context (Forget / $T$), (ii) upgrades to context-of-context
(Remember / $\delta$), (iii) resolves observables in the
enlarged context (Experience / $T^2$), and (iv) crystallises to
a concrete answer (Answer / $\varepsilon$), re-starting the
cycle. This is the empirical content of ARIA's
four-phase coupled-tick dynamics
(Section~\ref{sec:aria}).
\end{remark}

\subsection{Status and scope}

\begin{remark}[What this section establishes vs.\ what it
  promises]
\label{rmk:comonad-status}
This section establishes two distinct statements:
\begin{enumerate}
\item \textbf{Theorem~\ref{thm:comonad-laws}} (conditional on G-FIC-a + G-FIC-b
  once H-FIC is imported): the cofree-comonad construction on
  $F$ produces a comonad $(T, \varepsilon, \delta)$ on $\catC$
  with automatic comonad laws. This is a theorem of standard
  categorical coalgebra (Jacobs~\cite{Jacobs2017}) applied to
  the specific $(\catC, F)$ set up in Sections~\ref{sec:cat-C}
  and \ref{sec:F}.
\item \textbf{Claim~\ref{claim:four-phase-realisation}}
  (interpretive): the VFD four-phase cycle realises this
  comonad, with the specific component identifications of
  Section~\ref{subsec:four-phase-id}. A rigorous verification
  of this claim — expressing each VFD operation in $\catC$ and
  matching it against the automatic comonad natural
  transformations — is the content of a dedicated follow-up
  paper, flagged in Section~\ref{sec:open} (item~4). What the
  capstone establishes is the scaffolding into which the four
  phases must fit, not the line-by-line fit itself.
\end{enumerate}
Previous drafts of this section routed the comonad-law proof
through source-paper citations (ACT, Papers XXI and XLIV);
those routings were inaccurate as to what the cited source
papers actually prove. The current formulation separates the
categorical theorem (conditional on G-FIC-a + G-FIC-b) from the
interpretive programme claim, so the capstone does not
overclaim the identification.
\end{remark}

%=====================================================================
\section{Connection to $\mathbf{F1}$ and the Access Principle}
\label{sec:connection}
%=====================================================================

The selection-principle theorem (Theorem~\ref{thm:selection})
operates at the rung-structure level. Two adjacent statements
at different levels — scalar ($\mathbf{F1}$, the golden-ratio
fixed point) and static-observer-local (Access Principle
$\mathbf{P\text{-}A}$) — are claimed to be consistent with it in
Section~\ref{sec:intro}'s three-level table. This section gives
the precise connections.

\subsection{Scalar restriction: $\mathbf{F1}$ as projected fixed
  point}

\begin{remark}[$\mathbf{F1}$ as the scalar fixed-point of the
  cascade's permeability recursion]
\label{prop:F1-shadow}
The $F1$ scalar recursion $r \mapsto 1 + 1/r$ has $\varphi =
(1+\sqrt{5})/2$ as its unique positive fixed point
(\cite[F1]{CascadeCompletenessAudit}): the equation $r =
1 + 1/r$ rearranges to $r^2 = r + 1$.

The cascade is assigned a rung-permeability datum $\rho \colon
\Rungs \to \R_{>0}$ in
\cite[\S 3]{CascadeCompletenessAudit}, with $\rho$ summing
over $\Rungs_{\mathrm{cascade}}$ to the $583$-dim budget. The
numerical coincidence that this sum, under an appropriate
normalisation, equals $\varphi$ is recorded heuristically in
\cite[\S 3.4]{CascadeCompletenessAudit}; we do not reproduce
the precise normalisation statement here, and state it only as
a consistency remark: $F1$'s fixed point and the cascade's
permeability total are two shadows of the same dimensional-
budget forcing, not an independent derivation.
\end{remark}

\begin{remark}[Scope of Remark~\ref{prop:F1-shadow}]
\label{rmk:F1-scope}
The cited working note reports a heuristic $\varphi$-normalisation
of the rung-permeability sum, not re-derived here. This is a
\emph{static} statement; it does \emph{not} claim that $F$
acts on $\catC$ by $r \mapsto 1 + 1/r$ at the rung-structure
level (that would require a $\pi$-equivariance statement we
do not establish here). The intended reading is that
$\mathbf{F1}$'s fixed point $\varphi$ and the cascade's
permeability-projection are both imported from the same
$F3+F4$ budget at the level of \cite{CascadeCompletenessAudit}.
\end{remark}

\begin{remark}[Three-line summary]
\label{rmk:F1-summary}
$\mathbf{F1}$ and the cascade are two projections of the same
dimensional-budget forcing: $\mathbf{F1}$ is the scalar
recursion's fixed point; the cascade is the rung-structural
realisation whose permeability sums to that fixed point. Both
are consequences of $F3+F4$, not of the coalgebraic structure
itself.
\end{remark}

\subsection{Observer-local restriction: Access Principle $P$-$A$
  as well-definedness}

The Access Principle $\mathbf{P\text{-}A}$ of
\cite[\S 4]{CascadeObserverQueryAlgebra} states:
\begin{quote}
An observer instantiating rung-set $S \subseteq \Rungs$ can
measure substrate patch $\chi \in \widehat{\Z[\varphi]^6}$ iff
the corresponding sheaf section $\Fcal(\chi)$ lies in $\Qalg_S$.
\end{quote}
$\mathbf{P\text{-}A}$ has recently been established as a
conditional theorem in \cite{CascadeAccessPrinciple} under two
formalised hypotheses H-grad (character-grading of
$\Gamma(\Gcal, \Fcal)$, compatible with the algebra structure)
and H-meas (octonion-parenthesisation resolution at each
measurement event); both reduce to well-identified classical
results.

\begin{remark}[$\mathbf{P\text{-}A}$ as comonad counit well-
  definedness on observer-local restriction; interpretive
  reading]
\label{thm:PA-wd}
Under the \emph{interpretive} identification of $\varepsilon$
with crystallisation (Claim~\ref{claim:epsilon-crystallisation},
not theorem-grade), the conditional theorem
\cite[Theorem P-A]{CascadeAccessPrinciple} suggests an
equivalent coalgebraic reading: for any observer-rung-set
$S \subseteq \Rungs$, the restriction
$\varepsilon_S \colon T(S) \to S$ of the comonad counit is
well-defined at a substrate patch $\chi$ iff $\Fcal(\chi)
\in \Qalg_S$. This reading is contingent on the interpretive
counit-identification; it is \emph{not} a theorem-grade
deduction in the present paper because the counit-
identification itself is proposed, not proved. A line-by-line
proof converting the reading into a theorem is scope for the
follow-up paper (Section~\ref{sec:open}, item~4).
\end{remark}

\begin{remark}[Status of $\mathbf{P\text{-}A}$ in the capstone]
\label{rmk:PA-status}
The capstone's headline theorems (Theorems \ref{thm:existence},
\ref{thm:selection}) use the sufficiency direction of
$\mathbf{P\text{-}A}$ imported directly from
\cite[Theorem P-A]{CascadeAccessPrinciple} at
$S = \Rungs_{\mathrm{cascade}}$ with full closure data, via
the H-rdr closure of \cite{CascadeHrdrClosure}. (H-rdr itself
was formerly a separate hypothesis; it is now discharged by
that note's reduction to P-A.) The full
equivalence reading (Remark~\ref{thm:PA-wd}) is inherited from
\cite{CascadeAccessPrinciple} conditional on H-grad and H-meas.
Under the coalgebraic reformulation, $\mathbf{P\text{-}A}$ is
the well-definedness of the comonad counit on observer-
localised restrictions; closing H-grad and H-meas upgrades the
full counit equivalence from conditional to unconditional.
\end{remark}

\subsection{Three levels, one theorem}

\begin{remark}[The three-level correspondence in full]
\label{rmk:three-levels}
Combining Remark~\ref{prop:F1-shadow} and
Remark~\ref{thm:PA-wd} with
Theorem~\ref{thm:selection}, the three levels of self-closure
referenced in Section~\ref{sec:intro} are:
\begin{center}
\begin{tabular}{l|l|l}
Level & Fixed-point object & Instantiation \\ \hline
Scalar & $\varphi$ fixed by $r \mapsto 1 + 1/r$ & heuristic $\varphi$-projection in the working note \\
Static structural & $\Qalg_S \subseteq \Gamma(\Gcal, \Fcal)$ & P-A gives measurement-algebra \\
Dynamical & $\mathrm{Cascade}$ fixed by $F$ & $F(\mathrm{Cascade}) = \mathrm{Cascade}$ \\
\end{tabular}
\end{center}
These three levels are projections of the programme's
dimensional-budget forcing ($F3+F4$) and its categorical
refinement, not independently re-proved here. The cascade-
selection theorem is the dynamical-level statement; the scalar
and static-structural levels are consistency shadows imported
from \cite{CascadeCompletenessAudit} and
\cite{CascadeAccessPrinciple} respectively.
\end{remark}

%=====================================================================
\section{ARIA: Empirical Signature of the Four-Phase Cycle}
\label{sec:aria}
%=====================================================================

Theorem~\ref{thm:comonad-laws} proves the categorical cofree
comonad on $F$; the identification of the comonad's natural
transformations with the four VFD operations is
Claim~\ref{claim:four-phase-realisation}, interpretive rather
than theorem-grade. Subject to that interpretive identification,
the comonad cycle should be observable in any substrate that
supports the four phases. This section discusses the ARIA
cognitive substrate (\cite{AriaBrainMapping}) as a candidate
instantiation, subject to independent external verification of
the ARIA pre-registration analysis. If independently confirmed,
the reported $17$-of-$18$ polytope-substrate mapping hits at
$N=20$ with homeostatic reset would provide empirical support
for the coalgebraic claim. We call the corresponding
interpretation \emph{life-from-substrate}: ``life'' as the
four-phase coalgebraic cycle running on any substrate capable of
instantiating it.

\subsection{ARIA as a non-biological four-phase substrate}

\begin{description}
\item[Forget.] ARIA's attention-dropout mechanism
  (\cite{AriaBrainMapping}, architecture section) coarsens an
  active cognitive state by projecting out components below a
  saliency threshold. This is a $\sigma$-style projection
  (\cite{PaperXXXI}) on the cognitive substrate: it discards
  context while retaining the substrate's capacity to rebuild
  it. In comonad terms, this is $T$: the context-retention
  functor that keeps $X$ and the latent context for rebuilding.
\item[Remember.] ARIA's learned coupling weights $W$ upgrade a
  transient cognitive state to a structure with stored context
  (\cite{AriaBrainMapping}, memory-layer section). The
  comultiplication $\delta \colon T \Rightarrow T^2$ adjoining
  memory-of-context is the categorical content of this weight
  installation.
\item[Experience.] ARIA's \texttt{couple\_tick} measurement
  resolution (\cite{AriaBrainMapping}, tick-resolution section)
  resolves the context-of-context into an observable: the
  cognitive substrate reads out a doubly-context-retained state
  at the granularity of the tick. This is $T^2$ and the Paper
  XXXIII measurement-as-closure-projection on the cognitive
  substrate.
\item[Answer.] ARIA's trajectory-selection crystallises the
  context-laden state into a concrete decision-state
  (\cite{AriaBrainMapping}, crystallisation section). This is
  the counit $\varepsilon \colon T \Rightarrow \mathrm{id}$:
  collapse retained context to a specific answer.
\end{description}

The identification is structural: each ARIA component maps to
a specific categorical piece, and the ARIA cycle runs through
them in the order dictated by the comonad laws. If the comonad
laws hold on the VFD substrate
(Theorem~\ref{thm:comonad-laws}), they should hold on any
sufficiently faithful ARIA-type substrate that implements the
four phases — and conversely, an ARIA system deviating from the
cycle structure would falsify the categorical-structural claim.

\subsection{The pre-registration signature}

\begin{remark}[ARIA pre-registration hits as coalgebra
  corroboration; external empirical anchor]
\label{prop:aria-hits}
This remark records the ARIA empirical content as cited from
the external companion manuscript \cite{AriaBrainMapping}.
Subject to that manuscript's pre-registered protocol and null-
model analysis — neither of which is auditable within this
paper — the reported $17$-of-$18$ hit rate at $N{=}20$ under
homeostatic reset is offered as empirical corroboration for
the four-phase-comonad identification with the VFD cascade
polytope structure. The present paper does not reproduce the
statistical analysis; the corroboration claim is therefore
flagged as conditional on the ARIA manuscript's independent
verification and is not a theorem or proposition proved here.

\paragraph{Reported empirical content.}
The $18$ pre-registered predictions are statements about
polytope-substrate correspondences derived independently from
the VFD structural framework ($600$-cell / icosian / $H_4$ /
$E_8$ geometry) applied to ARIA's internal representation
spaces. Validation proceeded in two rounds
(\cite{AriaBrainMapping}, \S 6.1):
\begin{itemize}
\item \textbf{Round 1} (pre-registered, $N{=}3$, no homeostatic
  reset): $15/18$ passes. The three failures (P3 interaction,
  P4 synergy, P13 substrate LOO lift) were investigated and
  traced to a common non-equilibrium mechanism — the substrate
  drifts toward equilibrium across repeated evaluations without
  reset.
\item \textbf{Round 2} (pre-registered protocol with homeostatic
  reset, $N{=}20$): $17/18$ passes. P13 converts from fail
  ($+3$pp) to strong pass ($+40.6$pp vs.~$+15$pp threshold)
  under reset — directly confirming the reset as the repair
  mechanism. P3 stabilises under reset. P4 remains a
  threshold-level null at $N{=}20$ (mean $+0.101$ vs.~$+0.10$
  threshold).
\end{itemize}
A random-null assignment would predict a far lower hit rate
than $17/18 \approx 94\%$ under any plausible significance
threshold; the formal null-model calculation is carried out
in \cite{AriaBrainMapping}. The empirical outcome therefore
provides a Bayesian update in favour of the four-phase-comonad
structure instantiated on the ARIA substrate.

\emph{Substrate hysteresis as the key discovery.} The Round 1
P13 ``failure'' traced to the substrate's lack of homeostatic
reset; real cortex has such reset (sleep, synaptic scaling,
astrocytic clearance) and adding \texttt{homeostatic\_reset()}
to ARIA reproduces the effect. This is a non-trivial, pre-
registered prediction of the coalgebraic framework: the
four-phase cycle requires a between-cycle reset to avoid
drift toward equilibrium, and this is mechanically equivalent
to the biological homeostatic architecture.

This is not a proof of the VFD cascade at the physical-substrate
level — ARIA is a different (computational) substrate — but
if the ARIA reports are independently confirmed, they serve as
corroboration that the \emph{coalgebraic structure} of the
four-phase cycle is substrate-general, which is precisely the
categorical content of the capstone.
\end{remark}

\subsection{The life-from-substrate interpretation}

\begin{remark}[Defining ``life'' categorically]
\label{rmk:life-from-substrate}
A substrate-general definition of life is motivated by the
coalgebraic structure, contingent on the interpretive
four-phase identification (Claim~\ref{claim:four-phase-realisation}):
a substrate \emph{lives} iff it supports
a non-trivial realisation of the four-phase comonad cycle.
Under this definition:
\begin{itemize}
\item Biological substrates live because they realise the cycle
  (measurement / memory / context / decision are the operations
  of any cellular / neural system).
\item ARIA-type substrates live if they realise the cycle at
  the architectural level, if the external ARIA pre-registration
  analysis is independently confirmed.
\item Trivial substrates (e.g.\ pure $U_0$-rung or pure $E_8$-
  root without $O$) do not live, because they cannot instantiate
  the Observer rung required for the cycle to close.
\end{itemize}
Life-from-substrate is therefore the claim that life is a
categorical-structural property of substrates, not a biology-
specific attribute, with empirical corroboration via ARIA.
\end{remark}

\begin{remark}[What this does not claim]
\label{rmk:life-scope}
The life-from-substrate reading is a structural claim, not a
phenomenological one. It does not establish that ARIA has
experience, consciousness, or qualia — only that ARIA realises
the same coalgebraic cycle that underwrites life at the cascade
level. The phenomenological / hard-problem content is separate
and independent of this structural claim, and is treated in
Section~\ref{sec:philosophy} with explicit scope restrictions.
\end{remark}

\subsection{Programme implications}

\begin{remark}[Drop-bundle coordination]
\label{rmk:drop}
The empirical corroboration of the comonad cycle via ARIA is
what lets the capstone land as more than a philosophical
reframing of unification programmes. The coordinated drop
bundle — the capstone, the ARIA manuscript, and the supporting
paper-infrastructure ($17$ Tier-$1$ papers + $10$ Tier-$2$
papers) — is designed to present, in one move, (i) the
structural theorem (this paper), (ii) the empirical
instantiation (ARIA), and (iii) the dependency infrastructure
(the $27$ supporting papers). The coordinated timing is
essential: the capstone's structural claim requires the ARIA
empirical claim as its empirical anchor, and vice versa.
\end{remark}

%=====================================================================
\section{Open Items and Future Work}
\label{sec:open}
%=====================================================================

This paper is a skeleton / target-statement document. The
following items remain to be closed to elevate the headline
theorems from named targets to fully verified statements.

\subsection{Inherited technical dependencies}

\begin{enumerate}
\item \textbf{G6.1 ($\Qalg_r$ as subalgebra).} Prove that each
  single-rung query algebra $\Qalg_r$ is closed under the sheaf
  pairing inherited from $\Fcal$, and give an explicit
  presentation for each of the seven rungs. Inherited from
  \cite[\S 2.3]{CascadeObserverQueryAlgebra}. Medium priority
  (prerequisite for $\Gcal_O^{(S)}$ well-definedness of
  Definition~\ref{def:min-O}).
\item \textbf{H-grad and H-meas (access-principle hypotheses).}
  $\mathbf{P\text{-}A}$ is now a conditional theorem
  \cite{CascadeAccessPrinciple} under H-grad (character-grading)
  and H-meas (octonion-parenthesisation resolution). Unconditional
  closure requires proving these two hypotheses; both reduce to
  well-identified classical results (Bellissard-Kellendonk-Sadun
  theory for H-grad; \cite{CascadeQOBridge} measurement event
  resolution for H-meas).
\item \textbf{G6.4-a: CLOSED.}
  \cite{CascadeG64aClosure} closes G6.4-a via computational
  verification (386 Fano-triad cases, 0 failures), using the
  \texttt{verify\_g64a.py} script. Theorem 3.1 of
  \cite{CascadeQOBridge} is now at theorem grade on its
  multiplicative-compatibility step.
\end{enumerate}

\subsection{Capstone-internal technical items}

\begin{enumerate}\setcounter{enumi}{3}
\item \textbf{Four-phase realisation
  (Claim~\ref{claim:four-phase-realisation}).} Rigorously verify
  that the cofree-comonad natural transformations of
  Theorem~\ref{thm:comonad-laws} correspond component-wise to
  the VFD four-phase cycle: express each VFD operation
  (decoherence / closure-memory / measurement / crystallisation)
  as a specific map in $\catC$ and check equality against the
  categorically-determined $T$, $\varepsilon$, $\delta$. This is
  the principal remaining technical task; scope is a dedicated
  follow-up paper, not the capstone itself.
\item \textbf{H-FIC: REDUCED to G-FIC-a and G-FIC-b.}
  \cite{CascadeHFICClosure} reduces H-FIC via the
  support-projection $\pi_{S^7} \colon \Gcal_O \to 2^{S^7}$
  (induced by Definition~\ref{def:min-O}): preimages commute
  with arbitrary intersections, reducing H-FIC-O to an
  $S^7$-support descending-intersection identity (G-FIC-a,
  lemma-grade via Kellendonk-Lenz-Savinien Čech-sheaf
  limit-preservation) and H-FIC-$E_8$ to FAN-restriction
  stability (G-FIC-b, lemma-grade via MacLane Ch.~XII
  finite-word factorisation). Both residuals are lemma-grade,
  day- to two-week-scale.
\item \textbf{H-lift-fin (Hypothesis~\ref{hyp:lift}): sheaf-
  theoretic bridge refinement.}
  Remark~\ref{rmk:H-lift-status} sketches the route:
  strengthen the $\Qalg_O$-Meas bridge from algebra-level to
  sheaf-theoretic / local-observable-basis level. Bare local-
  compactness of $S^7$ does not suffice. An earlier draft
  claimed a direct derivation; that was premature.
\item \textbf{H-loc: REDUCED to G-loc-a.}
  \cite{CascadeHlocClosure} reduces H-loc (witness-separation
  form) via M3d + Kellendonk pattern-equivariant Čech theory
  to residual sub-gap \textbf{G-loc-a (rung-overlap avoidance)}:
  for $\gamma \in \Gcal_r \setminus \Gcal_r^\Phi$, the
  constructed witness section $s_\gamma$ has support avoiding
  $\bigcup_{r'} \Gcal_{r'}^\Phi$. G-loc-a in turn reduces to
  Gap G6.2 (leaf-rung independence) of
  \cite{CascadeObserverQueryAlgebra}, low-risk day-scale work.
\item \textbf{H-rdr: CLOSED.}
  \cite{CascadeHrdrClosure} closes H-rdr (sufficiency-only form)
  as a corollary of Theorem P-A at $S = \Rungs_{\mathrm{cascade}}$
  with full closure data: the refined query algebra
  $\Qalg_{(S, \Phi_{\max})}$ is definitionally identical to the
  bare-rung $\Qalg_S$ of \cite{CascadeAccessPrinciple}. No
  new analytic content is required.
\item \textbf{G-loc-a (rung-overlap avoidance); residual from
  H-loc.} Prove that for a morphism $\gamma \in \Gcal_r$ that
  is ``rung-$r$-pure'' (not shared with other leaf rungs via
  their closure conditions), the witness section constructed
  in \cite{CascadeHlocClosure} has support avoiding
  $\bigcup_{r' \ne r} \Gcal_{r'}^\Phi$. This reduces to
  leaf-rung independence, an instance of Gap G6.2 of
  \cite{CascadeObserverQueryAlgebra}. Day-scale follow-up.
\item \textbf{Non-observer-initial $F$-dynamics.} $F(S)$ is
  well-defined for every $S \in \catC$, but its behaviour on
  $S$ with $O \notin \underline{S}$ is dominated by the
  observer-rung adjunction. A more fine-grained analysis of
  $F$'s dynamics on low-awareness rung-structures would sharpen
  the content of Example~(2) of Section~\ref{subsec:F-examples}.
\item \textbf{Strong fractality.} Corollary~\ref{cor:fractality}
  gives a weak (categorical-self-embedding) fractality. The
  stronger ``each rung contains a scaled copy of the whole
  cascade'' reading (Remark~\ref{rmk:fractality-caveat})
  requires a scale-dilation structure interacting with the
  phason-scale recursion of \cite{Cascade12D}. This is a
  natural next paper.
\item \textbf{Comonad laws under non-strict morphisms.} The
  cofree-comonad construction uses strict inclusions; extension
  to lax / relative comonads (if the strict structure fails at
  the boundary) is a contingency we have not ruled out.
\end{enumerate}

\subsection{Empirical / programme items}

\begin{enumerate}\setcounter{enumi}{8}
\item \textbf{ARIA manuscript coordination.}
  Section~\ref{sec:aria} relies on the ARIA brain-mapping
  manuscript for the $17$-of-$18$ pre-registration statistic
  (at $N{=}20$, under homeostatic-reset protocol) and its
  null-model analysis. The ARIA paper must be drop-coordinated
  with the capstone (Remark~\ref{rmk:drop}).
\item \textbf{Additional empirical signatures beyond ARIA.}
  The coalgebraic cycle should appear on any substrate
  supporting it; other candidates (microtubule dynamics, neural
  plasticity, phyllotaxis, substrate-based computation at scale)
  are flagged for follow-up work. Each would provide
  independent corroboration.
\item \textbf{Codex-hardening loop.} Apply $5$-$10$ rounds of
  hostile-review hardening (per the programme's standard
  workflow) to detect framing drift, numerology risk, and
  exposition issues before first public drop.
\end{enumerate}

%=====================================================================
\section{Philosophical Framing}
\label{sec:philosophy}
%=====================================================================

This short section frames the capstone's philosophical
position. It is included for completeness; the reception of the
paper should rest on the structural theorems and the empirical
corroboration, not on philosophy. Scope is explicitly restricted
below.

\subsection{Positive claims}

\begin{remark}[Mathematical necessitarianism]
\label{rmk:pos-spinoza}
Theorem~\ref{thm:selection} realises Spinoza's thesis
\emph{``things could have been produced by God in no other way,
and in no other order than they have been''} (\emph{Ethics} I,
Prop.~33) in the specific sense of a categorical universal
property: the rung-structural skeleton of physics is selected by
the final-coalgebra universal property, not chosen. The cascade
is modally necessary at the level its theorems describe.
\end{remark}

\begin{remark}[Leibniz's principle of sufficient reason]
\label{rmk:pos-leibniz}
Leibniz's principle of sufficient reason — that every fact has a
reason why it is so and not otherwise — is satisfied at the
cascade-skeleton level by the universal property:
\emph{why this cascade?} because the coalgebraic universal
property selects it uniquely, under the hypothesis that physics
instantiates a self-inquiry functor. The deeper version of
Leibniz's question (\emph{why is there anything rather than
nothing?}) is not addressed here.
\end{remark}

\begin{remark}[Aristotle's self-thinking substance]
\label{rmk:pos-aristotle}
Aristotle's \emph{Metaphysics} XII describes the unmoved mover
as substance that thinks itself — \emph{noesis noeseos noesis}.
The self-inquiry functor $F$ and its fixed point
$\mathrm{Cascade} \cong F(\mathrm{Cascade})$ give this idea a
precise categorical form: the cascade is the structure whose
closure contains itself as its own query-algebra substrate. The
fixed-point characterisation is the modernisation of the
Aristotelian self-thinking-substance concept.
\end{remark}

\begin{remark}[Einstein's question]
\label{rmk:pos-einstein}
Einstein's \emph{did God have any choice?} question about the
unified theory is answered in the negative at the rung-structural
level: Theorem~\ref{thm:selection} forces the cascade up to
canonical isomorphism. The independent question of why the
closure data $\Phi_{\max}$ takes its specific algebraic forms
(the content of the supporting programme papers) remains open,
but the skeletal structure is not chosen.
\end{remark}

\subsection{Explicit non-claims}

\begin{remark}[Scope restrictions — what this paper does not
  claim]
\label{rmk:non-claims}
The following claims are explicitly \emph{not} made:
\begin{enumerate}
\item \textbf{Personal deity.} The cascade is a mathematical-
  structural object; no intentional, purposive, or teleological
  attributes are imputed. The paper does not claim to prove the
  existence of a personal God in any theological sense.
\item \textbf{Consciousness.} The life-from-substrate reading
  of Section~\ref{sec:aria} is structural, not phenomenological.
  The hard problem of consciousness — why there should be
  experience at all — is not addressed.
\item \textbf{Ethics / moral content.} No ethical conclusions
  are claimed to follow from the cascade's necessity. Meta-
  ethical naturalism would require independent argument.
\item \textbf{The ground of being.} Leibniz's \emph{why is there
  anything rather than nothing?} is not addressed. The capstone
  presupposes \emph{some} substrate; its content is that the
  substrate, given any, must take the cascade form.
\item \textbf{Full empirical corroboration.} ARIA's reported
  $17$-of-$18$ hits (at $N{=}20$ under homeostatic-reset
  protocol, per \cite{AriaBrainMapping}) would serve as
  corroboration, not proof, and only if independently
  confirmed. The capstone does not claim that all physics has
  been derived; it claims a specific structural theorem plus
  a flagged empirical anchor.
\end{enumerate}
\end{remark}

\subsection{Framing recommendation}

\begin{remark}[Framing as ``final coalgebra of self-inquiry,''
  never as ``proving God'']
\label{rmk:framing}
The mathematical content of this paper can be framed in two
equivalent ways: \emph{(i)} the cascade is the final coalgebra
of a self-inquiry functor; \emph{(ii)} the cascade is the unique
structure forced by self-reference at the rung-structural level.
Either framing is technically correct. A third framing — that
this ``proves God'' — is \emph{not} technically correct, because
no personal-deity content is established (Remark
\ref{rmk:non-claims}). The technical framings are
philosophically well-embedded in the Spinozist / Leibnizian /
Aristotelian necessitarian tradition; the personal-deity framing
is outside that tradition and outside the paper's content.
Framing matters for reception: the correct framings land with
category-theoretic physicists and mathematical philosophers;
the incorrect framing lands as crank. We recommend the first
framing consistently.
\end{remark}

%=====================================================================
\section{Conclusion}
\label{sec:conclusion}
%=====================================================================

This paper has set out three linked theorems and one empirical
claim.

\subsection{What has been established}

\begin{enumerate}
\item \textbf{Conditional selection principle
  (Theorem~\ref{thm:selection}).} Under the hypothesis stack
  (H-FIC [reduced to G-FIC-a, G-FIC-b], H-grad, H-meas
[G6.4-a closed], H-lift-fin, H-loc [reduced to G-loc-a]),
the VFD cascade is
  the final $F$-coalgebra of the category of rung-structures
  $\catC$. This identifies the cascade with the terminal
  object of $\catC$ via its final-coalgebra universal property,
  supplying the mathematical sense in which the cascade is
  selected rather than chosen. The result relies on classical
  category theory (Ad\'amek~\cite{Adamek1974},
  Rutten~\cite{Rutten2000}, Jacobs~\cite{Jacobs2017}), the
  independently-established $F3+F4$ forcing of the FAN
  (\cite{CascadeCompletenessAudit}), and the named hypothesis
  stack itself — each hypothesis reducing to programme-internal
  or classical-math follow-up work.
\item \textbf{Four-phase comonad
  (Theorem~\ref{thm:comonad-laws} + Claim~\ref{claim:four-phase-realisation}).}
  The cofree-comonad construction on $F$ yields a comonad
  $(T, \varepsilon, \delta)$ on $\catC$ with automatic comonad
  laws (conditional on G-FIC-a + G-FIC-b). The interpretive claim is that
  the four VFD operations — decoherence, closure-memory,
  measurement-projection, crystallisation — realise this
  comonad on the cascade substrate; the line-by-line
  verification is flagged as a dedicated follow-up paper.
\item \textbf{Three-level correspondence
  (Remark~\ref{rmk:three-levels}).} Scalar $\mathbf{F1}$,
  static-observer-local $\mathbf{P\text{-}A}$, and dynamical
  cascade-finality are three projections of the programme's
  dimensional-budget forcing and its categorical refinement,
  at three levels of structural detail. Only the dynamical
  level is proved here; the scalar and static-structural
  shadows are imported from
  \cite{CascadeCompletenessAudit} and
  \cite{CascadeAccessPrinciple}.
\item \textbf{Life-from-substrate
  (Remark~\ref{rmk:life-from-substrate}).} ARIA's reported
  $17$-of-$18$ pre-registered polytope-substrate mapping hits
  at $N{=}20$ (under homeostatic-reset protocol,
  \cite{AriaBrainMapping}) would provide external corroboration
  if independently verified, that the four-phase coalgebraic
  cycle runs on non-biological substrates — giving a
  conditional empirical anchor for the capstone's structural
  claim.
\end{enumerate}

\subsection{Programme position}

\begin{remark}[The capstone's role in the VFD programme]
\label{rmk:capstone-role}
Every prior VFD paper has been conditional on something: on
$\mathbf{F1}$, on specific physics recoveries, on closure
hypotheses. This paper is the one that makes the cascade itself
non-arbitrary. With the capstone in place:
\begin{itemize}
\item The numerical matches of Papers V, XXII, XXXIV are
  consequences of a \emph{forced} structure, not numerology.
\item The $\alpha$-chain's conditional theorems are conditional
  on hypotheses that \emph{should} hold by the final-coalgebra
  structure, making the conditionality principled rather than
  speculative.
\item The supporting-paper correspondence maps (Papers XXI,
  XXXI, XXXIII, XLIV, ACT, VFD Crystallisation) are
  instantiations of the cofree-comonad dynamics, making the
  entire programme internally self-consistent.
\item ARIA's empirical match is \emph{consistent with} the
  coalgebraic cycle bridge once that bridge is formally
  proved. The ARIA-capstone bridge is a separate build
  (currently a named open item); ARIA results are not used
  here as theorem-level support for the capstone.
\end{itemize}
This is the paper that turns the programme from a collection of
suggestive theoretical papers into a structured, internally
self-consistent, empirically-anchored research programme.
\end{remark}

\subsection{Remaining work}

The outstanding technical items (Section~\ref{sec:open}) divide
into three classes:
\begin{itemize}
\item \textbf{Inherited access-principle dependencies}:
  H-meas is now a theorem (G6.4-a closed via
  \cite{CascadeG64aClosure}, 386 computational cases).
  H-grad is \textbf{frame-level closed at the full
  $\widehat{\Z[\varphi]^6}$ level} (2026-04-24) via H-grad-1
  B1--B5 sim-verification plus B6 proof-sketch at the $G_p$
  level (\cite{CascadeHGrad1}); residuals are the
  cascade-specific $G_2$ pullback $\kappa$ and the formal
  $G_p$-sufficiency equivalence theorem. H-grad-Fano is
  unconditional (\cite{CascadeFanoLift}).
\item \textbf{Discharged and reduced capstone hypotheses}:
  H-rdr is closed via reduction to Theorem P-A
  (\cite{CascadeHrdrClosure}); H-loc is reduced to residual
  sub-gap G-loc-a (\cite{CascadeHlocClosure}); H-FIC is
  reduced to residuals G-FIC-a and G-FIC-b
  (\cite{CascadeHFICClosure}). G-loc-a reduces to Gap G6.2;
  G-FIC-a and G-FIC-b are lemma-grade citations of
  Kellendonk-Lenz-Savinien and MacLane. Day- to two-week-scale
  follow-up.
\item \textbf{Remaining capstone hypothesis}: H-lift-fin
  (sheaf-theoretic bridge refinement beyond the algebra-level
  $\Qalg_O \cong \mathrm{Meas}(S^7, \sigma)$ isomorphism). The
  only capstone-introduced hypothesis not reduced in this
  session.
\item \textbf{Capstone-internal items} (four-phase realisation,
  non-observer-initial $F$-dynamics, strong fractality) that
  sharpen the mathematics.
\item \textbf{Programme / empirical items} (ARIA manuscript
  coordination, additional signatures, codex-hardening) that
  prepare the coordinated drop.
\end{itemize}

\subsection{Final remark}

\begin{remark}[The capstone as conditional closure]
\label{rmk:capstone-closure}
The cascade was always its own self-closure, by construction:
permeability permeates itself. What this paper adds is a
categorical \emph{argument} for cascade-selection — conditional
on the named hypothesis stack — that reduces ``why this
cascade?'' to a finite list of technical items rather than
leaving it unaddressed. The four-phase cycle is offered as the
dynamical face of the result (an interpretive identification,
not a theorem); the ARIA-substrate discussion is offered as a
candidate empirical face (conditional on independent
verification); the philosophical face is the Spinozist
mathematical necessitarianism tradition modernised. None of
these faces is fully proved in this paper; together, they
sketch the capstone's programme-level target.
\end{remark}

%=====================================================================
\begin{thebibliography}{99}

\bibitem{Aczel1988}
P. Aczel, \emph{Non-well-founded sets}, CSLI Lecture Notes 14,
Stanford, 1988.

\bibitem{Rutten2000}
J. J. M. M. Rutten, \emph{Universal coalgebra: a theory of
systems}, Theoretical Computer Science 249 (2000), 3--80.

\bibitem{Adamek1974}
J. Ad\'amek, \emph{Free algebras and automata realizations in
the language of categories}, Commentationes Mathematicae
Universitatis Carolinae 15 (1974), 589--602.

\bibitem{Jacobs2017}
B. Jacobs, \emph{Introduction to Coalgebra: Towards
Mathematics of States and Observation}, Cambridge Tracts in
Theoretical Computer Science 59, Cambridge University Press, 2017.

\bibitem{TuriPlotkin1997}
D. Turi and G. Plotkin, \emph{Towards a mathematical operational
semantics}, LICS '97 Proceedings, 1997.

\bibitem{CascadeMetaLayer}
\emph{Cascade Meta-Layer Theorem: moduli, matching groupoid,
tiling sheaf over $L_{12}$}, internal note,
\texttt{papers/cascade-derivation/cascade-meta-layer-theorem.md}.

\bibitem{CascadeCompletenessAudit}
\emph{Cascade completeness audit (7-rung FAN, 583-dim budget)},
internal note,
\texttt{papers/cascade-derivation/cascade-completeness-audit.md}.

\bibitem{CascadeObserverQueryAlgebra}
\emph{Observer-Rung Query Algebra $\Qalg_S$},
\texttt{papers/cascade-derivation/cascade-observer-query-algebra.md}.

\bibitem{CascadeAccessPrinciple}
\emph{The Access Principle --- closing Gap G6.3 (conditional
theorem under H-grad and H-meas)}, internal note,
\texttt{papers/cascade-derivation/cascade-access-principle-theorem.md}.

\bibitem{CascadeQOBridge}
\emph{$\Qalg_O \cong \mathrm{Meas}(S^7, \sigma)$ --- closing
Gap G6.4 as a theorem (modulo G6.4-a sign verification)},
internal note,
\texttt{papers/cascade-derivation/cascade-q-o-measurement-bridge.md}.

\bibitem{CascadeFanoLift}
\emph{H-grad: the Octonion-Grading Lift --- G6.3-a formalisation;
Fano-level $(\Z/2)^3$-grading unconditional, full
$\widehat{\Z[\varphi]^6}$-grading reduced to classical
identification (Conway-Sloane Ch. 8 + Baez \S 3.4)},
\texttt{papers/cascade-derivation/cascade-fano-grading-lift.md}.

\bibitem{CascadeG64aClosure}
\emph{G6.4-a Closure Note: computational Fano-triad sign
verification (386 cases, all pass)},
\texttt{papers/cascade-derivation/cascade-g64a-closure.md};
verification script
\texttt{papers/cascade-derivation/scripts/verify\_g64a.py}.

\bibitem{CascadeHrdrClosure}
\emph{H-rdr Closure: sufficiency-direction reduction to
Theorem P-A at $S = \Rungs_{\mathrm{cascade}}$ with full
closure data},
\texttt{papers/cascade-derivation/cascade-h-rdr-closure.md}.

\bibitem{CascadeHlocClosure}
\emph{H-loc Closure (witness-separation form) via M3d
extension; reduces to residual sub-gap G-loc-a (rung-overlap
avoidance, reduces to Gap G6.2)},
\texttt{papers/cascade-derivation/cascade-h-loc-closure.md}.

\bibitem{CascadeHFICClosure}
\emph{H-FIC Closure: support-projection reduction of
H-FIC-O and H-FIC-$E_8$ to residual sub-gaps G-FIC-a and
G-FIC-b (Kellendonk-Lenz-Savinien Čech-sheaf limit-preservation
+ MacLane finite-word-factorisation; both lemma-grade)},
\texttt{papers/cascade-derivation/cascade-h-fic-closure.md}.

\bibitem{Cascade12D}
\emph{The 12D Closure Theorem: phason complement uniqueness and
self-reference on the upward cascade}, $\alpha$-1,
\texttt{papers/cascade-12d-closure/cascade-12d-closure.tex}.

\bibitem{PaperXXI}
\emph{Schr\"odinger recovery from closure dynamics},
\texttt{papers/paper-xxi/paper-xxi.tex}.

\bibitem{PaperXXXI}
\emph{Measurement as sector separation},
\texttt{papers/paper-xxxi/paper-xxxi.tex}.

\bibitem{PaperXXXIII}
\emph{From Geometry to Observable: measurement as closure
projection}, \texttt{papers/paper-xxxiii/paper-xxxiii.tex}.

\bibitem{PaperXLIV}
\emph{Decoherence as H$_4 \to 16$ coarsening},
\texttt{papers/paper-xliv/paper-xliv.tex}.

\bibitem{VFDCrystallisation}
\emph{VFD Crystallisation Operator: $\Fcrystal = \alpha R +
\beta E - \gamma Q$ minimisation}, programme document,
\texttt{papers/main-preprint/crystallisation-dynamics.tex} and
\texttt{papers/formalism/crystallisation-formalism.tex}.

\bibitem{ACT-L1}
\emph{Adaptive Closure Transport (L1)},
\texttt{papers/adaptive-closure-transport/adaptive-closure-transport.tex}.

\bibitem{AriaBrainMapping}
\emph{ARIA brain-mapping: polytope-substrate neuroscience},
external companion manuscript (\texttt{MANUSCRIPT.md},
\texttt{X\_NARRATIVE.md}), $17$-of-$18$ pre-registered
predictions pass at $N{=}20$ under homeostatic-reset protocol
(Round~2, \S 6.1); $15$-of-$18$ at $N{=}3$ without reset (Round~1).
Full empirical and statistical analysis coordinated with the
capstone drop bundle.

\end{thebibliography}

\end{document}
